\documentclass[../main.tex]{subfiles}
\begin{document}
%==========================================TIÊU ĐỀ TẬP========================================
\begin{table}[h]
  \begin{adjustwidth}{-1cm}{}
    \begin{tabular}{>{\centering\arraybackslash}p{6cm}|>{\raggedright\arraybackslash}p{12.5cm}}
      \multirow{5}{6cm}
      % Cột bên trái: ÉPISODE và số tập
      {\\[-40pt]
        \flaregothic\fontsize{20pt}{20pt}\selectfont \centering
        \color{eptype}ÉPISODE\color{black}\\
        \burbank\fontsize{72pt}{72pt}\selectfont
        \color{epnum}95\color{black}\\[6pt]
        \cafeteria\fontsize{20pt}{20pt}\selectfont\color{epnum}(7-11)\color{black}
        \\[18pt]
        % \flaregothic\fontsize{14pt}{14pt}\selectfont
        % \color{epattr}BIENVENUE, \uline{\color{Banpire} Mel} !
        \flaregothic\fontsize{18pt}{18pt}\selectfont
        \color{epattr}FIN DE TOME !\\
        \color{black}
      }
       &
      % Dòng thứ nhất: tên tập tiếng Nhật
      {\fotskip\fontsize{18pt}{18pt}\selectfont ショッピングモールで\ruby{暴}{ぼう}\ruby{走}{そう}したい！
      }    \\[6pt]
      % Dòng thứ hai: tên tập tiếng Anh
       & {\cafeteria\fontsize{18pt}{18pt}\selectfont Cleaning Up!}                       \\[4pt]
      % Dòng thứ ba: tên tập tiếng Việt
       & {\vnmsans\fontsize{18pt}{18pt}\selectfont Cùng dọn dẹp và trang trí căn phòng!} \\[8pt]
      % Dòng thứ tư: Nhân vật xuất hiện
       & {\caviar{\fontsize{14pt}{14pt}\selectfont Nhân vật: \emp{kuon1} \emp{ryuuhou1} \emp{sora} \emp{roboco2} \emp{miko2} \emp{mio} \emp{korone} \emp{noel} \emp{polka} \emp{lamy} \emp{botan}}
        }
      \\[8pt]
      % Dòng cuối cùng: Thời gian diễn ra của tập (thường sẽ là ngày phát hành)
       & {\continuum\fontsize{16pt}{16pt}\selectfont Ngày phát hành: 28/02/2021}         \\[18pt]
    \end{tabular}
  \end{adjustwidth}
\end{table}
%=========================================TRUYỆN CHÍNH========================================
\color{black}

\txtbx[2]
{{\color{vol} \uline{Tóm tắt:}} Cuối tháng Hai, Kuon cùng Roboco và Noel dọn sạch đống thùng bìa trong văn phòng, khiến Mio, Korone và Miko ngạc nhiên. Roboco và Noel được nhận tem ``Cảm ơn vì đã giúp!'' thay vì tiền lương. Ryuuhou cũng tham gia và khoe tem với anh trai. Cả nhóm đi mua đồ trang trí: Miko mua cục đá Moai, Botan mua đèn hình lựu đạn, Polka mua đống hề nhồi bông, Lamy mua lọ hoa thủy tinh, còn Ryuuhou mua khung ảnh, kệ treo và đèn dây. Họ còn chiến đấu với một con gián và xịt thuốc khắp phòng, rồi cùng nhau trang trí văn phòng. Cuối cùng, Sora tặng tem cho mọi người.
}

Ngày cuối cùng của tháng Hai, mùa đông đang dần lùi lại để nhường chỗ cho mùa xuân đang tới. Dù thời tiết vẫn còn hơi lạnh, những tia nắng ấm đã bắt đầu xuất hiện, báo hiệu một mùa mới sắp bắt đầu. Trong không khí trong lành đó, mọi người đều cảm thấy phấn chấn, tràn đầy năng lượng để bắt đầu những kế hoạch mới.

Và một trong những việc ấy chính là dọn dẹp, quét sạch toàn bộ đống thùng carton trong văn phòng, việc đáng lẽ phải làm ngay từ ngày thành lập công ty, cách đây đã hai năm rồi mà mãi mới đến lúc thực hiện.

Bộ ba Mio, Korone và Miko, cùng ba thành viên Thế hệ thứ Năm là Polka, Botan và Lamy bước vào phòng. So với văn phòng cũ đầy tối tăm và lộn xộn trước đây vốn chất đầy những thùng các-tông phía sau chiếc ghế sofa màu xanh lá cùng bàn uống trà bừa bộn và vách ngăn gỗ, căn phòng lúc này đã hoàn toàn thay đổi: toàn bộ không gian ngập tràn ánh sáng trắng rực rỡ, rộng rãi và sạch sẽ đến kinh ngạc, thay thế ghế sofa bằng những dãy bàn làm việc và tủ thấp màu trắng gọn gàng, chỉ còn lại chiếc chổi quét nhà cùng một quả bóng đá nằm ở góc phòng.

\img{7-13a}

Nàng Khuyển thốt lên đầy phấn khích, đôi mắt cô sáng rực như một đứa trẻ vừa thấy căn phòng mới: ``Văn phòng hôm nay sao rộng thế này!''

Nàng Sói đen cũng gật gù đồng tình, trong khi Lamy đưa tay che miệng, mắt tròn xoe ngạc nhiên nhìn quanh: ``Ủa? Ngày nào chỗ này cũng chật cứng thùng hàng mà nhỉ?''

Nàng Cáo xiếc nhìn quanh một lượt, rồi nghiêng đầu thắc mắc với vẻ bối rối: ``Hửm? Chẳng nhẽ hôm nay mình vào nhầm phòng rồi sao?''

Bất chợt, một giọng nói cất lên từ đằng sau Korone. Miko giật mình, thét lên đầy cảnh giác: ``\textbf{Ai đó!?}''

Cô quay đầu về phía giọng nói, và thấy bộ ba Roboco, Noel cùng Kuon bước ra từ góc phòng, bên cạnh còn có em gái Ryuuhou, cũng đang tươi cười. Roboco phẩy tay, đáp lại: ``Bọn chị đây đã dọn hết đống hộp trong phòng rồi!''

Korone híp mắt, khoanh tay nghi hoặc như một thám tử đang điều tra vụ án: ``Là ba siêu lính ngự lâm đã làm việc này sao!?''

\img{7-13b}

Mio cũng nhướng mày, nhìn về phía cậu Tổng: ``Kuon-senpai cũng dọn cùng hai người đó à?''

Kuon thở dài, khoanh tay nhìn cả đám với gương mặt của người đã quá quen với việc phải thu dọn mớ hỗn độn: ``Đáng lẽ chúng ta nên làm vậy từ hai năm trước\dots\ Nhưng mấy đứa quậy quá, cộng với cấp trên cũng phàn nàn nhiều, nên anh và Ryuuhou với hai đứa PON này mới quyết định đi dọn.''

Cậu chỉ hai ngón cái vào hai cô gái, vẻ mặt ngán ngẩm. Ngay lập tức, Roboco và Noel cùng đồng thanh phản đối: ``Bọn em không có PON!!''

Bên cạnh, cô bé mắt hai màu cười khúc khích, giơ xấp tem nhỏ lấp lánh lên khoe với anh trai: ``Onii-san, em cũng hăng hái dọn dẹp lắm nhé! Noel-san và Roboco-san đã phụ giúp khuân vác rồi, em chỉ có việc quét dọn với lau nhà thôi!''

Kuon toát mồ hôi hột, nhìn em gái vừa bất lực vừa ngạc nhiên: ``Sao em cũng dính vào vụ này thế?''

Ryuuhou hớn hở đáp, giọng như một đứa trẻ đang khoe chiến tích: ``Em thấy mọi người dọn dẹp nên xin vào phụ giúp thôi!''

Kuon nhướn mày, thở dài bất lực: ``À\dots\ nhà mình thì bảo dọn một góc phòng cũng không thèm, sang đây dọn văn phòng thì mày hăng hái nhất nhỉ? Cái kiểu đúng là `việc nhà thì nhác, việc chú bác thì siêng' ha.''

Ngay lập tức, cô em gái phồng má phản đối, giậm nhẹ chân không đồng tình: ``Em có lười đâu mà anh nói! Hôm trước em còn lau nhà ở nhà mà!''

Câu chuyện kéo đến bên phía Yêu tinh tuyết đứng ngoài lắng nghe cũng chen vào hỏi Kuon: ``Kuon-senpai, nếu không gọi mấy người bạn thì ạm định gọi ai dọn thật hả?''

Cậu Tổng tặc lưỡi, rồi quay sang em gái với vẻ bất lực: ``Kệ đi, anh không biết nữa. Anh thấy hai đứa khỏe nên nhờ, chứ không thì anh cũng gọi mấy bạn qua dọn là xong.''

Botan khoanh tay, bật cười trước màn đối đáp của họ: ``Chẳng trách mà ba người dọn nhanh thế.''

Kuon cười trừ, đáp lại như người đã quen với mọi điều bất ngờ: ``Thì anh nghĩ không chỉ mọi người ngạc nhiên đâu, đến cả A-san cũng vậy.''

Cậu hướng mắt về phía nàng Đeo kính đang ngồi trên ghế xoay. Polka hí hửng chạy tới, ngồi xổm xuống như một phóng viên điều tra: ``Vậy vụ này có liên quan đến A-chan đúng không!?''

Cậu Tổng bật cười, cầm một cái micro không khí và giả vờ phỏng vấn với giọng dẫn chương trình chuyên nghiệp: ``Và sau đây, nhân chứng Yuujin A-san có đôi lời muốn nói. Chị cảm thấy thế nào về hai người ấy ạ, thưa chị A-san?''

Nàng Đeo kính thoáng cười, đáp lại một cách điềm tĩnh, đầy hài lòng: ``Ồ, họ giúp bọn tôi rất nhiều! Đống thùng bìa đó chỉ mất năm phút là sạch bách luôn!''

``Ể!? Nhanh đến vậy sao?'' Lamy ngạc nhiên trầm trồ, mắt mở to không thể tin nổi. Botan thì bật cười với phản ứng của bạn mình, thán phục: ``Vậy mà em cứ tưởng phải mất ít nhất nửa tiếng chứ.''

``Vâng, và đó là cảm nghĩ của chị A-san. Xin mời quay trở lại trường quay với Omaru-san.''

Bên kia, Mio vỗ nhẹ vào đầu Noel, hào hứng khen ngợi như một người chị đang thưởng công em gái: ``Chị mừng lắm, hy vọng hai người sẽ được trả đủ lương cho công sức dọn dẹp này đấy!''

Noel ngay lập tức khoanh tay, ngẩng cao đầu tự hào đáp lại đầy kiêu căng: ``Đương nhiên rồi, chả có gì phải bàn!''

Roboco chỉ cười khúc khích, vẻ mặt có vẻ đang giấu giếm điều gì đó, đôi mắt long lánh lộ rõ vẻ bí ẩn.

``Ơ\dots\ thế là cả ba người đều được nhận tiền thưởng sao!?'' Miko thốt lên đầy kinh ngạc, giọng lộ rõ sự ghen tị vì sợ mình bị bỏ quên. Bản thân cô cũng ao ước được thưởng, ánh mắt đầy vẻ thèm muốn.

Cùng lúc đó, Korone đang dùng ngón tay tạo thành một cái ống kính, đặt trước mắt rồi lần lượt xoay chuyển tay này qua tay khác như đang quay phim, gương mặt nghiêm trọng chẳng khác nào một nhiếp ảnh gia chuyên nghiệp.

Kuon lắc đầu liên tục, vẫy tay phủ nhận ngay: ``Không đâu. Thực ra chỉ có hai người được nhận lương thôi, bởi vì\dots''

Cậu chỉ tay về phía Roboco và Noel. Lập tức, hai cô gái cười đầy bí ẩn, vừa giơ cao xấp tem nhỏ lấp lánh vừa trông vô cùng tự hào.

Noel nở nụ cười tươi rói, giọng phấn khích như một đứa trẻ đang khoe thành tích mới đạt được: ``Bọn em nhận được vô số tem `Cảm ơn vì đã giúp đỡ!' á!''

Roboco cũng liên tục gật gù, thêm vào đầy hãnh diện: ``Nhiều đến mức dán kín cả quyển sổ tay của em luôn ấy!''

Ryuuhou nhìn chằm chằm vào xấp tem trên tay hai người, lẩm bẩm đầy hớn hở: ``Tem đẹp thật! Em cũng phải dán chúng vào sổ tay của mình mới được!''

Cả căn phòng bỗng chốc im lặng. Tất cả ánh mắt đều đổ dồn về phía ba con người đang cầm tem, không ai tin vào điều vừa xảy ra. Polka chớp mắt nhìn hai người, rồi quay sang Botan cùng Lamy, giọng nói đầy bối rối.

``Chúng ta có đang lạc vào một thế giới khác hay không vậy?'' nàng Cáo sa mạc toát mồ hôi hột, vừa sợ hãi với xấp tem vô tri đó vừa liên tục lắc đầu như thể đang cố gắng tỉnh táo.

``Ừm\dots\ có lẽ đây là một vũ trụ song song rồi,'' nàng Sư tử cũng tỏ ra nghi ngờ, hai mắt híp lại đầy hoài nghi.

``Không lẽ\dots\ những cái tem đó có thể quy đổi ra tiền mặt thật sao?'' nàng Yêu tinh tuyết mở to mắt kinh ngạc, buột ra câu hỏi đầy bất ngờ.

Kuon chỉ cười nhạt, khoanh tay giải thích như người đã quá quen với những điều kỳ quặc xảy ra hàng ngày: ``Nó giống mấy cái 'phiếu bé ngoan' ở nơi anh làm trước đây thôi, chỉ có giá trị tinh thần thôi, chứ không thể đổi trác được gì đâu.''

\img{7-13c}

Mio lập tức phản đối gay gắt, giọng đầy bất bình: ``Mấy chị là nhi đồng thối mũi à!?''

Miko cũng trợn tròn mắt, than thở: ``Hai người bị người ta lợi dụng rồi mà còn vui vẻ gì, ne!''

Kuon thấy hai tiền bối phản ứng thái quá, liền bình thản đáp lại như đang nói về một điều hiển nhiên: ``Anh thấy chuyện này cũng bình thường mà.''

``Nhận tem thay lương mà anh còn nói là bình thường ư!?'' nàng Vu nữ không thể tin được vào sự vô tư của cậu Tổng, giọng nói đầy kinh ngạc. Kuon thở dài, gương mặt như người đã chứng kiến quá nhiều chuyện lạ trong đời, thanh minh: ``Không, cái việc dọn dẹp văn phòng ấy. Sớm muộn gì cũng phải làm thôi. Ai lại thèm mấy cái tem trẻ con đó đâu, anh chỉ nói việc dọn dẹp là bình thường thôi.''

Lúc này, nàng Khuyển cũng lên tiếng phán xét đầy nghiêm túc, giọng chẳng khác nào quan tòa đang tuyên án trước tòa: ``Nhưng kiểu trả công thế này không phải là vi phạm luật lao động à?''

Chẳng ngờ, ngay sau khi cô nói xong, Roboco và Noel đi tới chỗ A-chan và áp giải cô đi mất, như hai cảnh sát đang bắt tội phạm. Họ thực sự đã tin vào lời đó, gương mặt đầy nghiêm túc và quyết tâm.

Vừa bị lôi đi, nàng Đeo kính vừa ngơ ngác nhìn mọi người, nói trong bối rối: ``Ơ? Vi phạm á? Từ từ, đừng có còng chị chứ\dots''

Lamy che miệng khúc khích cười, trong khi Botan khoanh tay lắc đầu với vẻ hết thuốc chữa, còn Kuon thì ngán ngẩm buông một câu: ``Anh nghĩ là mấy đứa hiểu nhầm gì đó rồi\dots\ Động não suy luận đi chứ.''

Ryuuhou đứng bên cạnh, vẫn cầm xấp tem trên tay, nhìn cảnh tượng trước mắt vừa buồn cười vừa bối rối: ``Onii-san, các chị ấy đang làm gì thế? Chị A-chan bị bắt thật à?''

Kuon chỉ thở dài, đưa tay xoa đầu em gái: ``Hy vọng là không phải bắt vào tù thật.''

\ct

Sau một hồi mất công giải thích đầu đuôi mọi chuyện cho cả Roboco và Noel, cùng với nhân viên bảo vệ tòa nhà - người đến kiểm tra sau khi có người báo cáo về vụ ``bắt giữ trái phép'' trong văn phòng, cuối cùng A-chan cũng được thả tự do. Hai cô ``ngốc'' sau đó được mọi người đón về nhà, trước khi kịp thực hiện cái kế hoạch vô lý là đổi những chiếc tem lấy bánh ngọt từ máy bán hàng tự động: một ý tưởng chỉ những tâm hồn quá ngây thơ mới có thể nghĩ ra. Ryuuhou cũng đứng bên cạnh anh mình, cười khúc khích nhìn bộ đôi vừa bị ``tống cổ'' về.

Lamy thở phào nhẹ nhõm, giọng cô như tiếng thở dài của sự bình yên vừa trở lại: ``May mà không ai bị kéo lên phường thật nhỉ.''

Botan khoanh tay, bật cười với giọng đầy thích thú: ``Ừ, dù gì thì việc nhận tem thay lương cũng không phải là tội hình sự mà.''

Polka gật gù, rồi chợt nảy ra một ý tưởng lạ với vẻ mặt khám phá ra chân lý mới: ``Thế nếu mình mang mấy tem đó đổi lấy hiện vật thì được không? Kiểu như đổi lấy vé vào công viên giải trí chẳng hạn?'' Ryuuhou đứng bên cạnh cũng hùa theo, nhón chân gật đầu liên tục như đồng tình hoàn toàn với đề xuất của cô nàng Cáo xiếc.

Kuon híp mắt nhìn cô, cười méo mặt với giọng cố gắng không bật cười lên: ``Anh không nghĩ là A-san có hệ thống đổi thưởng hoành tráng như thế đâu\dots\ Cả mấy đứa nữa, không biết mấy cái phiếu đó là để làm gì thật sao?''

Nàng Đeo kính A-san đẩy nhẹ cặp kính, lắc đầu với giọng của một nhân viên văn phòng từ chối yêu cầu vô lý: ``Ừ, xin lỗi nhé Omaru-san. Chị đâu phải là trung tâm đổi quà gì đâu.''

Cả đám cười rộ lên, tiếng cười giòn tan vang khắp căn phòng vừa được dọn dẹp sạch bách. Không khí thoáng đãng, dễ chịu hơn hẳn, như một làn gió mới thổi qua căn phòng từng chật chội năm nào. Ryuuhou cũng nhảy cẫng cười theo, chạy vòng quanh một vòng quanh căn phòng rộng rãi để cảm nhận không gian mới.

Với căn phòng bỗng dưng rộng thênh thang, ba người - Korone, Mio và Miko - tranh thủ ngồi trên ghế xoay để chạy vòng quanh, đua xem ai có thể lướt được xa nhất. Tiếng bánh xe lăn trên sàn nhà vang lên như bản nhạc vui nhộn của đám trẻ con. Nàng Khuyển có lẽ là người tận hưởng nhất, cười không ngớt mỗi khi lướt qua Kuon - người lúc này đang nằm dài trên ghế sô-pha ngả lưng, rồi còn cắt ngang trước mặt nàng Vu nữ làm cô giật mình co chân lên.

\img{7-13d}

Korone giơ hai tay lên cao như đang lướt ván, hô to với giọng phấn khích không che giấu: ``Rộng rãi quá! Nhưng mà\dots\ trông cũng trống trải quá nhỉ!''

Cậu Tổng lười biếng đáp lại, giọng cậu như lời thì thầm trong cơn mơ chiều: ``Nhìn lại phòng thì đúng là thế thật. Nhưng anh thấy, tối giản cũng tốt mà.'' Ryuuhou lúc này chạy lại ngồi bên cạnh anh, cũng ngả đầu vào vai Kuon mà gật gù đồng tình với lời của anh mình: ``Em cũng thấy đơn giản thế này là ngăn nắp và gọn gàng nhất.''

Miko chống cằm, nhăn mặt suy nghĩ như một nhà thiết kế đang đau đầu với bản vẽ chưa xong: ``Nhưng nếu thế thì nhìn sẽ nhàm chán ne! Bao quanh chỉ có bốn bức tường như thế này, chẳng có đồ trang trí gì cả\dots''

Mio cũng gật gù đồng tình với Korone và Miko, giọng cô như đang tìm kiếm một giải pháp hoàn hảo: ``Đúng là thế thật. Giá như ta có cách nào làm cho văn phòng này rực sáng hơn nhỉ\dots''

Bất chợt, cô nảy ra một ý tưởng hay, vỗ mạnh vào mặt bàn khiến cả phòng nghe thấy tiếng động: ``Muốn trang trí nó không? Ta có thể mua ít đồ về đặt trong phòng được đấy!''

Lập tức nàng Vu nữ vỗ tay cái ``bốp'', đồng tình ngay lập tức, nhoẻn miệng cười hào hứng: ``Ý hay nhe\~{}'' Còn Kuon thì lười nhác lên tiếng, giọng như người đã không còn quan tâm gì nhiều: ``Cũng được. Anh thì thế nào cũng ổn.'' Ryuuhou cũng nhảy cẫng lên vui sướng, kêu lớn ``Ông yeah! Đi mua đồ thôi!'' làm cả đám cười lớn.

Đúng lúc này, Polka chen ngang vào cuộc trò chuyện, giọng cô nghiêm chỉnh như đang đưa ra câu hỏi quan trọng nhất: ``Nhưng trang trí gì mới được? Chúng ta phải chọn chủ đề cho văn phòng chứ!''

Lamy gật gù, nghiêng đầu suy nghĩ như một nghệ sĩ đang tìm kiếm cảm hứng cho tác phẩm: ``Ừ nhỉ, nếu bày biện lung tung quá thì sẽ thành\dots\ rối mắt lắm.''

Botan thì đơn giản hơn, khoanh tay đề xuất với giọng người đã có kinh nghiệm trang trí: ``Cứ theo phong cách hiện đại, đơn giản mà ấm cúng là ổn rồi.''

Nghe xong ba cô em gái nói xong, Korone bỗng đứng bật dậy, chỉ tay về phía cửa với giọng đầy nhiệt huyết như một đội trưởng dẫn quân: ``Vậy thì, xuất phát nào!''

Korone quá phấn khích đến quên mất mình vẫn đang ngồi trên ghế xoay. Cô đạp mạnh một cú xuống sàn, khiến cái ghế phóng thẳng ra ngoài cửa, trượt trên hành lang như một viên đạn bay khỏi nòng, tốc độ tăng dần như một phương tiện mất kiểm soát, rồi đúng như dự đoán\dots\ *RẦM!*

Mio, Miko và Kuon lập tức chạy ra ngoài xem tình hình với vẻ mặt đầy lo lắng. Và đúng như họ sợ, nàng Khuyển bây giờ đang nằm sõng soài dưới chân cầu thang, tay chân dang rộng như một ngôi sao biển bị sóng đánh dạt vào bờ. Ryuuhou chạy theo sau, thấy cảnh tượng đó thì bật cười lớn đến mức ngồi sụp xuống.

Cả ba người lớn đồng loạt thở dài, nhìn nhau bó tay, rồi cùng thanh than vãn với giọng đầy bất lực: ``Sao lại phấn khích làm cái gì (ne)\dots''

Trong khi đó Polka vỗ tay cười khoái chí, thốt lên như một khán giả đang xem vở kịch hay nhất đời: ``Korone-senpai đúng là không bao giờ làm mình thất vọng!'', còn Botan và Lamy thì vừa che miệng cố nhịn cười, vừa không giấu nổi sự thích thú, đôi vai cả hai run lên vì cố kìm nén tiếng cười sắp bật ra.

\il

Từ một góc khuất phía sau khu vực máy tính, một đôi mắt xanh lá tròn xoe đang dòm ngó từng cử động của mọi người trong phòng. Cô đã trốn ở đây được một lúc rồi, ban đầu chỉ định thăm thú xem liệu có ai hay biết sự hiện diện của mình không, thế rồi chẳng mấy chốc lại bị cuốn vào câu chuyện đang diễn ra trước mắt, tựa như một khán giả vô tình xem được một bộ phim hài tự nhiên, không có kịch bản nào được sắp đặt trước.

Cậu Tổng đang đứng nói chuyện với ba người bạn cùng thế hệ của cô: một người tóc bạc, vẻ mặt lúc nào cũng điềm nhiên và thư thái; một người tóc xanh dịu dàng, luôn hết lòng chăm sóc mọi người xung quanh; và cuối cùng là một người tóc vàng năng động, chẳng bao giờ chịu ngồi yên một chỗ. Y như cô đã nhận định từ trước, họ lúc nào cũng tràn đầy niềm vui mỗi khi ở cạnh nhau.

Nhưng điều làm cô bối rối vô cùng là tại sao cậu ấy cũng có mặt ở đây?

Cô khẽ phồng má lên, tự trách bản thân sao vẫn chưa dám bước ra. Lẽ ra giờ này cô đã phải lao tới chào hỏi mọi người thật rộn ràng như mọi khi, thế mà đôi chân lại cứ như chôn chặt xuống nền nhà, chẳng thể nào nhấc lên nổi, như thể có cái gì đó đóng đinh chúng lại.

Thông thường, cô đâu phải kiểu người hay ngần ngại trước bất cứ điều gì. Thế mà hôm nay, chẳng hiểu tại sao tim lại cứ bồn chồn không yên, cảm giác y hệt một diễn viên nhí đang đứng sau cánh gà, chuẩn bị cho buổi diễn đầu tiên trên sân khấu lớn.

Cô lặng lẽ rụt đầu lại sau bức tường, nắm chặt hai bàn tay lại rồi hít thật sâu một hơi, cố gắng bình tĩnh lại.

``Chỉ cần đợi lúc thích hợp thôi\dots''

Dù thế nào đi nữa, cô cũng không thể trốn ở đây mãi được. Sớm muộn gì cô cũng phải bước ra, gia nhập cùng mọi người. Nhưng bây giờ, cô vẫn cần thêm chút ít thời gian để chuẩn bị tinh thần cho khoảnh khắc đó.

\il

\begin{center}
\uline{Trung tâm thương mại \AE ON - quận \ruby{Ryuuhen}{Long Biên}.}
\end{center}

Bảy người bước vào, ánh đèn rực rỡ phản chiếu trên nền gạch bóng loáng, tạo cảm giác hiện đại mà vẫn ấm áp. Cánh cửa tự động vừa mở là Korone đã hớn hở chạy nhảy tung tăng như một chú cún mới được thả xích, đôi mắt cô long lánh phấn khích như đứa trẻ lần đầu được đến công viên giải trí.

\img{7-13e}

``\textbf{Rộng quá!!}'' cô reo lên, rồi xoay tròn giữa sảnh chính, hai tay dang rộng như muốn ôm cả không gian vào lòng. Korone chạy lông nhông khắp trung tâm, cứ như đây là lần đầu tiên cô đặt chân đến nơi này, để lại phía sau dãy tiếng cười giòn tan của mọi người.

Mio đưa mắt nhìn quanh, ngạc nhiên hỏi: ``Đây là đâu vậy?''

``Ở JUSCO\footnote{Tiền thân của \AE ON trước khi được đổi tên vào năm 2001. Một số người Nhật vẫn dùng cái tên này nhằm nhớ lại thời xưa cũ.} đây cái gì cũng có đó!'' Miko đáp, hào hứng không kém, mắt đã dò xét các quầy hàng xung quanh với vẻ thèm thuồng của đứa trẻ đứng trước cửa hàng kẹo.

Nàng Sói đen toát mồ hôi hột, vội kéo lấy cánh tay Miko trước khi cô nàng nói thêm gì nữa: ``Miko-senpai, đó là \AE ON bây giờ rồi mà\dots''

Polka đứng bên cạnh cũng khúc khích cười, huých nhẹ cô bạn Yêu tinh tuyết bên cạnh để ám chỉ điều gì đó: ``Ối giời, Miko-obaasan nhà mình vẫn còn nhớ cái tên JUSCO cơ đấy, già thật rồi già thật rồi\~{}''

Kuon nhìn ba người với vẻ nửa cười nửa bất lực, giải thích như một hướng dẫn viên du lịch: ``Như Miko-chi nói đấy. Em cứ mừng tượng rằng đây là một siêu thị khổng lồ. Các gian hàng chia khu rõ ràng, và chắc chắn ở đây có cả đống thứ thú vị để mấy đứa khám phá.''

Ryuuhou, em gái của Kuon, đứng sát bên anh, mắt sáng rực: ``Onii-san, em chưa bao giờ đến trung tâm thương mại nào to thế này kể từ khi ta vào GO!\footnote{Tên mới của chuỗi siêu thị Big C từ tháng 12/2020.} từ năm năm trước! Nhìn kìa, có cả quán kem nữa!''

Người anh trai mỉm cười nhẹ nhàng nhắc nhở em gái: ``Ừ, nhưng nhớ là mình đến đây để mua đồ trang trí văn phòng đấy nhé. Đừng có mua cái gì không cần thiết đâu đấy.''

Mio gật đầu, vừa hiểu ra thì Polka đã vội kéo tay hai người bạn đồng thế hệ lại gần một quầy trưng bày, tiếng reo vui vang lên như đàn chim non tung cánh.

``Oa\~{} nhìn này, mấy con plushie này dễ thương ghê!'' nàng Cáo sa mạc tít mắt, nhấc lên một con gấu bông hình lạc đà Alpaca, mặt rạng rỡ thích thú.

Lamy mỉm cười dịu dàng, cố gắng giữ vẻ nghiêm chỉnh nhắc nhở: ``Ừm, ở đây bán đủ loại lắm. Nhưng đừng quên mục đích chính là mua đồ trang trí phòng nhé?''

Botan khoanh tay, nhìn quanh một lượt với vẻ điềm đạm của người đã có kế hoạch rõ ràng: ``Chúng ta cứ đi dạo một vòng trước đã. Dù sao thì cũng còn nhiều thời gian mà.''

Ryuuhou nhìn ba chị hậu bối đang hào hứng, khẽ thì thầm với anh trai: ``Onii-san, các chị ấy trông vui quá. Em có thể đi cùng các chị ấy một lát không?''

Kuon gật đầu đồng ý: ``Được, nhưng đừng đi xa quá nhé.''

Mọi người dần dần bị phân tán bởi các gian hàng hấp dẫn xung quanh. Cũng dễ hiểu thôi, đây là lần hiếm hoi - nếu không nói là lần đầu - họ được đặt chân vào nơi này, choáng ngợp trước ánh đèn lung linh và vô số hàng hóa đủ màu sắc. Trái ngược với hai anh em nhà Hikawa, nơi trung tâm thương mại và siêu thị không phải là một nơi gì đó quá xa lạ.

Chưa kịp trao đổi thêm được gì, nàng Korone đột nhiên đứng cứng ngắc trước cửa hàng thú cưng kế bên, đôi mắt tròn to lên, ngây ngất như vừa tình cờ tìm thấy thiên đường.

``Chú chó thật kìa!'' nàng reo lên, ngón tay chỉ vào chú cún nhỏ nằm trong lồng kính trưng bày, hai bàn tay vừa ấn chặt vào mặt kính như muốn được chạm vào bộ lông mềm mại của nó.

``Nhắc đến chó thì\dots'' cậu Tổng cất tiếng, rồi cùng với Mio, cả hai cùng dồn ánh nhìn vào một chú chó lớn hơn đang sủa váng ngay dưới chân; một con chó to bị lạc chủ, đang lang thang bơ vơ giữa dòng người qua lại tấp nập.

\img{7-13f}

Kuon nghiêng đầu, bối rối không thể hiểu nổi, cố gắng suy nghĩ để lý giải một sự việc vô lý: ``Sao lại có một con chó lạc giữa trung tâm thương mại thế này nhỉ?''

Ryuuhou vội chạy lại gần, cúi người xuống ngắm chú chó với vẻ mặt tròn mắt thích thú: ``Chú chó to ghê! Chắc là nó bị lạc mất chủ người rồi.''

Nàng Sói đen chỉ biết lắc đầu thở dài, rồi nhẹ nhàng bế chú chó lạc lên. Ở phía bên kia, Miko đột nhiên ngửi thấy một mùi thơm ngọt ngào thoảng bay trong không khí, mũi cô khẽ hít hà liên tục như một thợ săn đang dò theo mùi con mồi.

``Hửm? Mùi này là\dots''

Cô vừa dò tìm hướng mùi vừa hít hà liên tục, rồi chẳng mấy chốc lại hét lên phấn khích: ``\textbf{Người ta bán bánh cá taiyaki nè\~{}!}''

Chẳng chần chờ gì nữa, cô lập tức lao đi như một mũi tên bắn khỏi cung, phóng thẳng về phía quầy hàng, bỏ lại phía sau Mio đang đứng bất lực gọi theo: ``Miko-chi!?''

Kuon khẽ mỉm cười, lắc đầu ngao ngán, đã quá quen thuộc với sự bốc đồng, lém lỉnh của những cô bạn: ``Mấy đứa chắc cũng hơi hiếm khi có dịp đến trung tâm thương mại đông vui thế này nhỉ.''

Mio đưa tay lên xoa thái dương, trên mặt hiện rõ vẻ mệt mỏi của một người đang phải cố gắng trông nom một lũ trẻ con hiếu động: ``Đúng là thế thật\dots\ Thôi thì, mọi người nhớ giữ liên lạc với nhau là được rồi.''

``Ừ, ta cứ từ từ đi dạo một vòng khám phá. Chưa hết một giờ rưỡi mà, còn nhiều thời gian lắm,'' Kuon đáp, vừa nói vừa ra hiệu để mọi người có thể tự do tản ra khám phá trung tâm theo cách mình muốn.

\ct

Chẳng mấy chốc, trong dòng người qua lại tấp nập, cả nhóm tách ra từng mảng nhỏ, ai cũng bị hút vào những gian hàng đầy đồ vật hấp dẫn xung quanh. Nửa tiếng trôi qua nhanh hơn một cái chớp mắt, và khi Mio dừng chân trước một cửa hàng bán phụ kiện thú cưng, chợt nhận ra chỉ còn mình cô, Kuon và Ryuuhou đang đứng cùng nhau. Những người khác đã biến mất giữa biển người, chẳng để lại dấu vết nào ngoài tiếng cười xa xăm vảng đâu đó trong không trung.

Nàng Sói đen lôi điện thoại ra, nhắn vào nhóm chat chung với giọng điệu của người đang kiểm tra đám trẻ con tản mạn: ``Hai đứa đang ở đâu thế?''

Chỉ vài giây sau, nàng Khuyển nhắn lại trả lời ``Tại đây nè'' kèm một bức ảnh tự sướng cạnh mô hình người ngoài hành tinh và chiếc đĩa bay giả, trong ảnh đôi mắt cô lấp lánh rạng ngời.

\img{7-13g}

Mio trợn tròn mắt, giọng đầy hoảng sợ và bối rối: ``\textbf{Chỗ nào vậy trời!?}''

Ryuuhou ngó vào màn hình điện thoại của nàng Sói đen, cười khúc khích: ``Korone-san đang ở khu vui chơi đó! Nhìn chị ấy vui chưa kìa!''

Cậu Tổng cũng bật cười, giọng đầy bất lực trước sự hiếu động của nàng Khuyển: ``Đúng là khu vui chơi rồi. Làm sao mà em ấy chạy lăn tăn đến được nơi đó nhỉ?''

\il

Trong khi mọi người còn đang hào hứng khám phá khắp trung tâm thương mại - trừ cậu Tổng và cô em gái của cậu, những người đã quá quen thuộc với những nơi như thế này, và nàng Sói đen đang bối rối hoảng loạn - thì ở một góc xa, một cô gái tóc vàng ánh hồng đang núp sau biển quảng cáo, lén lút dõi theo từng bước chân của nhóm bạn. Dù đội mũ lưỡi trai và đeo kính râm để che giấu danh tính, trong mắt cô vẫn ánh lên niềm thích thú không thể giấu.

Khi ánh mắt lướt qua những người bạn cùng Thế hệ của mình, cô khẽ gật đầu, như thể vừa xác nhận được sự hiện diện của ba người bạn cùng thế hệ. Nhưng khi ánh nhìn dừng lại ở cậu Tổng, cô bỗng khựng lại, gò má tự nhiên ửng hồng lên một chút.

Cô nhanh chóng quay người đi, kéo vành mũ sụp xuống thêm một đoạn, tiếp tục quan sát từ phía sau một gian hàng bên cạnh.

Ryuuhou đang dạo bước lang thang gần đó, bỗng dừng chân. Cô bé nheo mắt nhìn về phía người lạ bí ẩn, rồi thì thầm nhỏ vào tai anh trai vừa bước đến: ``\hl{Onii-san, hình như có người đang theo dõi tụi mình kìa.}''

Cậu Tổng nhíu mày, ngước nhìn theo hướng em gái chỉ: ``Có người theo dõi thật à? Em chắc chứ?''

``Em chắc lắm. Chị đó cứ nhìn về phía nhóm mình hoài, nhất là những lúc anh đang nói chuyện với mọi người ấy.''

Cậu Tổng hướng mắt về phía tấm biển quảng cáo, nhưng người bí ẩn đã khéo léo lẩn mất. Cậu chỉ có thể lẩm bẩm trong hơi thở: ``\textit{Có gì đó đáng ngờ lắm đây\dots}''

Khi cô vẫn đang lén lút quan sát từ phía sau gian hàng, cô bất giác giật mình bởi một tiếng hét lớn vang vọng khắp cả trung tâm thương mại.

\il

Một tiếng hét vang lên từ phía Miko, người đang lao tới như một mũi tên bắn khỏi cung. Nghe giọng quen thuộc, Mio như tìm thấy vị cứu tinh, giọng đầy nhẹ nhõm: ``A! Miko-chi quay lại rồi kìa! Chị về--''

Câu nói bị cắt ngang khi cả nhóm nhìn thấy một chiếc bánh cá biết đi đang rượt đuổi nàng Vu nữ với tốc độ đáng kinh ngạc, y hệt một sinh vật kinh dị trong phim đang săn đuổi con mồi của mình.

\img{7-13h}

Theo phản xạ tự nhiên, cả nàng Sói đen và cậu Tổng đều lập tức chạy theo họ. Ryuuhou, đang đứng gần đó cũng vội vã bám sát phía sau, vừa chạy vừa thắc mắc đầy bối rối: ``Onii-san, tại sao Miko-san lại bị một cái bánh đuổi theo thế?''

``Anh cũng đâu có biết!'' cậu Tổng đáp lại, trong khi đôi chân vẫn không tự chủ di chuyển theo phản xạ.

Vừa chạy, nàng Sói đen vừa hớt hải hét lên với giọng đầy hoang mang: ``\textbf{Tại sao lại có chuyện kỳ quặc thế này xảy ra chứ!?}''

Nàng Vu nữ vừa thở hổn hển vừa kể lại đầu đuôi câu chuyện trong những hơi thở gấp gáp. Tóm lại: Miko tình cờ gặp một người mặc trang phục bánh cá taiyaki - chính là kẻ đang đuổi theo bây giờ - và tự dưng nảy ra hàng loạt câu hỏi triết lý vô nghĩa, như một nhà triết học đang thẩm vấn một hiện tượng siêu nhiên:

\begin{itemize}
\item ``Mày sinh ra trên đời để làm gì vậy?''
\item ``Tại sao mày lại bị nướng lên phục vụ người ta mỗi ngày chứ?''
\item ``Mày có bao giờ được nhận tem `Cảm ơn vì đã giúp đỡ!' hơm ne?''
\end{itemize}

Đương nhiên, ``người bánh cá'' đã tức điên lên, và thế là cuộc rượt đuổi điên rồ này xảy ra, biến một câu chuyện bình thường thành một câu chuyện cổ tích bị xoay chiều kỳ lạ.

Nàng Sói đen vừa chạy vừa hét lên đầy bất lực: ``\textbf{Chị đúng là ngốc nghếch! Mau xin lỗi người ta đi!}''

Ryuuhou vừa theo kịp mọi người, vừa cười khúc khích dù đang thở dốc: ``Miko-san đúng là có cách riêng để gây chuyện thật đấy, anh nhỉ?''

Cậu Tổng thì chỉ biết lắc đầu ngao ngán, giọng đầy bất lực: ``Bó tay.com luôn. Đang yên lành tự dưng đi chọc tức người ta làm gì cho ra nông nỗi này!?''

Cả ba người chạy vòng quanh tầng một của trung tâm thương mại suốt mười lăm phút, như một vở kịch hài đang được trình diễn trước sự ngạc nhiên của những người qua đường. Cuối cùng, họ chỉ có thể cắt đuôi được ``người cá takoyaki'' khi họ chạy vào một cửa hàng và đứng yên giả làm ma-nơ-canh, nín thở trong bóng tối.

Sau khi ``người cá'' chạy đi mất, ba người mới thở phào nhẹ nhõm, những hơi thở như vừa thoát khỏi một cơn ác mộng. Miko thì thầm với giọng còn run run: ``Ta an toàn chưa ne?''

``Hình như nó đi rồi đấy,'' cậu xuống bục, bỏ chiếc mũ vành xuống với vẻ mặt nhẹ nhõm.

Nàng Vu nữ thở phào nhẹ nhõm, giọng cô như một lời cầu nguyện: ``May quá rồi ne\dots\ Kami-sama (Chúa) phù hộ cho bọn mình rồi\dots''

``Em không nghĩ là Kami-sama sẽ tha tội cho chị đâu,'' nàng Sói đen liếc nhìn tiền bối, trưng ra khuôn mặt bất bình của một người đã chứng kiến quá nhiều sự bốc đồng.

``Cẩn thận bị nghiệp quật đến mức không ngóc đầu lên được đấy em ơi,'' cậu Tổng hùa thêm vào, giọng như một lời cảnh báo đầy hài hước.

``M-Mấy người cứ nói mấy cái quái gở gì thể không biết ne!?'' Miko đỏ mặt, vung vẩy hai cánh tay biện minh như một đứa trẻ bị bắt quả tang.

Ryuuhou đứng bên cạnh, cô bé cười khúc khích: ``Miko-san, lần sau đừng có hỏi mấy câu đó nữa nhé. Không thì lại bị đuổi tiếp đấy!''

Sau đó, cả ba quay trở lại với công việc chính: đi dạo một vòng quanh trung tâm thương mại, đồng thời tìm mua đồ để trang trí căn phòng.

Miko cuối cùng cũng tìm đúng quầy bán taiyaki và cô đã có một bữa ngon lành ở đó, thưởng thức từng miếng bánh với vẻ mặt hạnh phúc như một đứa trẻ.

Mio thì vô tình tìm thấy chủ nhân của một chú chó đi lạc, người đó cứ ríu rít cảm ơn cô mãi, còn cô nàng thì chỉ đáp lại với giọng khiêm tốn và một nụ cười đầy mãn nguyện: ``Dạ không, chị cứ khách sáo rồi\dots''

Korone vẫn trú ngụ ở khu vui chơi giải trí, nơi mà cô đang ``bay lơ lửng'' trong một mô hình trạm vũ trụ, tay cô dang rộng như một phi hành gia. Cô còn gặp gỡ và chơi đùa với hàng loạt mô hình quái vật ngoài hành tinh, tỏ ra phấn khích đến mức quên mất khái niệm thời gian, như một đứa trẻ trong công viên giải trí.

Botan không ngạc nhiên lắm trước sự rộng lớn của \AE ON nhưng lại thích thú với khu đồ công nghệ. Cô thử nghiệm hàng loạt thiết bị, từ tai nghe không dây, bàn phím cơ cho đến các máy chơi điện tử đời mới, đôi mắt cô sáng lên với sự thích thú của một người đam mê công nghệ.

Polka tìm thấy một cửa hàng bán đồ hóa trang, nơi cô thử một bộ trang phục trông giống như diễn viên xiếc, không quên tạo dáng trước gương với vẻ mặt hào hứng của một nghệ sĩ đang chuẩn bị cho buổi diễn.

Lamy thì khám phá khu mỹ phẩm và hương liệu, thử các loại nước hoa mới, thi thoảng còn nhắn tin cho nàng Cáo Sa mạc và nàng Sư tử trắng để hỏi xem loại nào phù hợp với mình nhất. Xong rồi, cô còn len lỏi hẳn sang cửa hàng bán rượu để tìm loại rượu phù hợp với mình, với vẻ mặt của một người sành sỏi.

Kuon và Ryuuhou là người duy nhất nghiêm túc trong việc mua đồ. Cậu, cùng với cô em gái và Mio sau đó, cẩn thận xem xét từng quầy hàng để tìm đồ trang trí phù hợp cho văn phòng, những bàn tay tỉ mỉ lựa chọn từng món đồ. Cuối cùng, họ mua được một số vật phẩm cần thiết và ngồi chờ bên ngoài khoảng nửa tiếng, thư giãn sau một buổi mua sắm mệt mỏi.

Cả bảy người họ đều tranh thủ chụp ảnh ở mỗi nơi họ ghé qua để làm kỷ niệm, những tiếng cười và tiếng máy ảnh vang lên không ngớt.

\img{7-13i}

Và trong những tấm hình đó, nếu để ý kỹ, sẽ thấy có một chi tiết hơi lạ.

Trong bức ảnh nàng Sói đen chụp cậu Tổng và cô em gái của cậu đang chọn đồ trang trí, ở phía xa có bóng dáng một cô gái đội mũ lưỡi trai và đeo kính râm, lấp ló sau một kệ hàng, như một bóng ma trong bức tranh.

Trong bức ảnh nàng Cáo Sa mạc chụp với bộ đồ xiếc, ở góc nhỏ bên trái, một mái tóc vàng nhạt lấp ló sau quầy tính tiền, như một ẩn số trong bức ảnh.

Trong bức ảnh nàng Yêu tinh tuyết chụp một quầy nước hoa với dòng chữ ``Hương thơm dịu dàng như một đóa hoa'', ngay sát mép tấm gương phản chiếu một hình bóng lờ mờ quen thuộc, như một linh hồn đang lảng vảng.

Ryuuhou, khi xem những tấm ảnh trên điện thoại của anh trai, bỗng nhiên nhíu mày: ``Onii-san, hình như có người cứ xuất hiện trong ảnh của chúng ta mãi vậy.''

Kuon nhìn em gái, rồi nhìn lại những tấm ảnh, một nụ cười nhẹ nở trên môi: ``Ừ, anh thấy rồi.''

Tuy nhiên, chỉ có ba người - nàng Cáo Sa mạc, nàng Yêu tinh tuyết và nàng Sư tử trắng - nhận ra điều này ngay lập tức, những ánh mắt họ trao nhau với một sự thấu hiểu không lời.

Và dĩ nhiên, cậu Tổng cũng biết. Cậu khẽ nhếch mép cười, giọng cậu như một lời thì thầm: ``Thế này là thế nào nhỉ\dots

\dots N███-chi?''

\ct

\begin{center}
\uline{Bốn giờ chiều.}
\end{center}

Bảy người cùng nhau bước chân trở về văn phòng, ai cũng tay xách nách mang đủ thứ đồ đạc, những chiếc túi đủ các màu sắc lủng lẳng trên tay với đủ vẻ mặt hớn hở khác nhau. Sau chuyến đi mua sắm vừa rồi, không khí cả nhóm vẫn còn tràn đầy năng lượng, tiếng bước chân xen lẫn tiếng cười nói ríu rít vang vọng suốt dọc hành lang.

Korone hớn hở reo lên, giọng đầy phấn khích như vừa thắng được cuộc thi nào, tức tốc đặt túi xuống sàn nhà: ``Về tới nhà rồi\~{}!!''

Miko cũng hào hứng nối lời ngay, giọng trong trẻo chẳng khác nào đứa trẻ vừa kết thúc chuyến dã ngoại vui vẻ: ``Bọn mình về tới nơi rồi đây!''

Lamy chỉ mỉm cười dịu dàng, nhẹ nhàng nhận xét: ``Chuyến đi hôm nay vui thật đấy nhỉ!''

Polka cũng gật gù xuýt xoa, vẻ mặt vẫn còn lưu lại sự ngạc nhiên sau khi khám phá một thế giới đồ sộ: ``Chẳng thể ngờ trong cái trung tâm thương mại đó lại có nhiều thứ đa dạng đến thế.''

Botan cũng đồng tình lia lịa, giọng tràn đầy sự thán phục: ``Nhiều đồ công nghệ hiện đại quá, nhìn đến mắt không rời.''

Mio thở dài một hơi, có chút tiếc nuối: ``Tiếc thật là cả nhóm không thể cùng nhau đi dạo và chọn đồ luôn. Mà Kuon-senpai, em không ngờ anh lại kén chọn kỹ lưỡng đến thế trong việc sắm đồ cho phòng đấy.''

Cậu Tổng cười đáp, giọng điềm nhiên như người đã có nhiều kinh nghiệm sắp xếp không gian: ``Phong cách trang trí cho căn phòng phải hợp với sở thích của tất cả mọi người chứ em. Làm sao mà cứ thấy gì thích là bừa bãi mua về được.''

Ở bên cạnh, Ryuuhou vẫn chưa động đến chiếc túi của mình, chỉ ôm nó trước ngực với nụ cười bí ẩn, rồi khẽ thì thầm vào tai anh trai: ``Onii-san, đợi lát nữa em mới khoe món đồ em mua nhé, giữ bí mật một chút!''

Cậu Tổng quay sang nhìn em gái với ánh mắt tò mò rộn lên: ``Mua gì mà giấu giếm ghê thế hả?'' Cô bé chỉ cười tinh nghịch, nháy mắt một cái: ``Anh cứ đợi một chút là biết liền!''

Mọi người trải hết tất cả đồ đạc xuống sàn, bắt đầu kiểm kê lại những món đồ vừa mua về, tiếng sột soạt của túi ni-lông và giấy gói vang lừng khắp căn phòng.

``Cả nhà mua được gì nào?'' Korone tò mò hỏi, đôi mắt cô long lánh tràn đầy háo hức.

Miko hăm hở đáp lời ngay: ``Lôi hết ra xem luôn ne!''

Mio nhìn đống túi đồ của mọi người, lộ rõ vẻ lo lắng, giọng cô như một lời cảnh báo nghiêm túc: ``Chị thì không có ý kiến gì đâu, nhưng mà mấy đứa có mua thứ gì kỳ quặc hay ngớ ngẩn không đấy? Đây là công ty idol đấy, nhớ rõ không?''

``Anh cá là mấy đứa chắc chắn đã về với thứ chẳng liên quan gì đến công chuyện đâu,'' Kuon tặc lưỡi, lắc đầu ngao ngán, chẳng hề tin tưởng gì tới họ.

Nghe hai người nghiêm túc nhất nghi ngờ, nàng Chó lập tức đáp lại với vẻ mặt vô cùng nghiêm chỉnh, giọng cô đầy tự tin như đang khẳng định một điều chắc chắn: ``Không có đâu! Hai người cứ tin bọn em đi.''

Nàng Vu nữ cũng gật đầu lia lịa, đồng thanh khẳng định: ``Đúng đó! Bọn chị đã bàn bạc thỏa thuận với nhau từ rồi mà!''

Nàng Sói đen dường như bị sự chắc chắn của hai cô nàng thuyết phục, cô ngạc nhiên, giọng có chút xúc động: ``Hai đứa à\dots\ Th-thật sự vậy hả?''

\dots Để rồi chẳng mấy chốc, cô mới nhận ra mình đã đặt niềm tin hoàn toàn sai chỗ.

Mio ôm một con mèo nhồi bông trong tay, Kuon thì xếp mấy chậu cây cảnh vừa mua vào một góc. Hai người là duy nhất mua về đúng những món đồ trang trí cho văn phòng. Còn Korone và Miko thì hoàn toàn ngược lại.

\img{7-13j}

Nàng Vu nữ đập mạnh xuống sàn một cục đá khổng lồ hình đầu người, to lớn, nặng nề và còn cao hơn cả đầu người, khiến cả nhóm đều giật mình. Tiếng động lúc đó vang lên như một tiếng sét đánh giữa trời,  nàng Sói đen trợn tròn mắt, hét lên với giọng đầy kinh ngạc và hoảng sợ: ``\textbf{Cái này lạ lùng, ngớ ngẩn quá mức rồi còn gì!}''

Cậu Tổng nhíu mày chặt chẽ, giọng cậu chẳng thể tin được vào mắt mình: ``Thế mà em vác được cái khối đá khổng lồ này lên đây được sao? Trông nó to đùng hết biết\dots''

``Xem chừng nó có vẻ nặng vài trăm cân chứ chẳng đùa,'' Ryuuhou nhìn cục đá đầy xét nét. ``Miko-san mang nó lên tới tận tầng này kiểu gì vậy\dots?''

Chủ nhân của cục đá chẳng vội trả lời, cô áp tai vào bề mặt đá lạnh lẽo, nói bằng giọng đầy bí ẩn như một nhà tiên tri đang lắng nghe lời sấm truyền từ trời cao: ``Hai người thấy không, em cảm nhận được nó đang gọi em đó ne\dots''

Lời giải thích đó chẳng thuyết phục được ai. Mio tức giận tháo chiếc kính râm đang đeo, cả người cô dường như bừng lên ngọn lửa phẫn nộ. Còn Kuon thì chỉ đứng đó chống hông, lắc đầu ngao ngán với giọng của một kẻ đã từ bỏ mọi hy vọng: ``Em bị ảo ma Canada thật rồi hay sao\dots''

Vừa dứt câu, Mio ném luôn chiếc kính râm của mình vào cục đá và hét lên trong cơn giận dữ: ``\textbf{Gọi cái đầu chị ấy!}''

*CỐP!*

Cục đá nghiêng ngả rồi đổ ập xuống sàn với một tiếng rầm kinh hoàng, làm cả căn phòng rung chuyển. May mắn là nó chưa làm hỏng cái gì, nhưng cũng đủ khiến cậu Tổng hoảng hốt, giọng cậu đầy lo lắng: ``Miko-chi\dots\ Làm ơn vứt cái cục đá này đi giùm anh đi.''

\img{7-13k}

Ba người còn lại - Botan, Lamy và Polka - nãy giờ vẫn đứng yên quan sát, cuối cùng cũng cất tiếng lên tiếng, khuôn mặt ai cũng vừa ngạc nhiên vừa thích thú trước cảnh tượng vừa xảy ra.

``E-Em cứ nghĩ chị sẽ mua thứ gì đó\dots\ nhẹ nhàng, dễ thương hơn chứ!?'' Lamy khẽ chạm vào bề mặt đá, ngơ ngác hỏi, ngón tay cô lướt nhẹ trên lớp đá lạnh lẽo.

Botan khoanh tay trước ngực, gật gù như thể đang đánh giá một tác phẩm nghệ thuật đỉnh cao, giọng cô đầy suy tư: ``Chà, nhìn cũng có nét độc đáo đấy chứ. Nhưng mà\dots\ để làm gì được bây giờ?''

Còn Polka thì bật cười khoái chí, vỗ tay bồm bộp như một khán giả đang xem vở kịch hay nhất đời: ``Hay quá Miko-senpai! Đây chính là phong cách độc đáo mà nhóm idol của mình cần! Lạ, độc và chẳng ai hiểu nổi!''

Kuon ôm trán, thở dài một hơi đầy mệt mỏi, giọng cậu của một người đã từ bỏ mọi hy vọng: ``Thêm ba đứa bất thường nữa rồi\dots''

Botan nhún vai, nở một nụ cười thong dong với giọng điệu như sắp khiến cả nhóm bất ngờ: ``Đừng có mất niềm tin sớm thế chứ, Kuon-senpai. Được rồi, giờ đến lượt bọn em nhỉ?''

Cô đặt chiếc hộp vuông vức xuống sàn, mở ra bên trong là một chiếc đèn bàn hình quả lựu đạn, màu xanh rêu với những đường nét sắc lẹm, trông chẳng khác nào món đồ chơi của ngừi lính.

Mio há hốc mồm, mặt cắt không còn giọt máu khi lùi lại mấy bước, giọng đầy hoảng sợ: ``\textbf{Botan!! Cái đó làm sao mà hợp với văn phòng của nhóm idol được chứ!!}''

Kuon ngồi xổm xuống, nhấc chiếc đèn lên để ngắm nghía, vẻ mặt giống hệt một nhà phê bình nghệ thuật đang xem xét tác phẩm: ``Mà công nhận trông cũng ngầu thật\dots\ nhưng mà có ai lại ngồi làm việc với một quả lựu đạn phát sáng trên bàn đâu chứ?''

Botan cười khẽ, bấm công tắc. Một ánh sáng dịu nhẹ tỏa ra từ chốt an toàn của chiếc đèn, không quá chói mà lại rất dễ chịu, như một tia nắng ấm xuyên qua màn đêm. Cô gật đầu hài lòng, giọng đầy tự hào về lựa chọn của mình: ``Trông có vẻ nguy hiểm thế thôi, mà ánh sáng rất ổn đúng không? Còn có cả chế độ chỉnh độ sáng nữa đấy.''

Miko gãi đầu, vẫn chưa thể hiểu nổi chuyện gì đang xảy ra, giọng đầy bối rối: ``Ừm\dots\ nghe cũng có lý đấy ne\dots\ nhưng mà vẫn lạ lắm!''

Kuon ngắm chiếc đèn thêm lúc lâu, vẻ mặt đầy phân vân: ``Ít nhất thì cái này cũng không gây nguy hiểm gì. Chỉ tội hình dạng nó dễ làm người khác hiểu lầm quá. Mang về trang trí phòng em thì hợp lý hơn đấy.''

Cậu đưa lại chiếc đèn cho nàng Sư tử, người đang vẫy đuôi ríu rít vì thích thú, y hệt một con thú cưng vừa tìm được món đồ chơi mới. Sau đó cậu quay sang người tiếp theo, giọng đầy chờ đợi: ``Rồi, còn em thì sao, Lamy-chan?''

Nàng Yêu tinh tuyết cẩn thận mở chiếc túi đồ của mình ra, động tác nhẹ nhàng như đang khám phá một kho báu bị lãng quên. Bên trong là bộ lọ hoa thủy tinh trong suốt, mỗi chiếc đều có họa tiết hoa tuyết lấp lánh, tinh xảo và thanh tao, giống hệt những bông tuyết vừa ngưng đọng trên mặt kính.

Mio thở phào nhẹ nhõm, như vừa thả được một tảng đá lớn khỏi ngực, giọng đầy sung sướng: ``Cuối cùng cũng có một người có gu thẩm mỹ bình thường!''

Lamy đỏ bừng mặt, lúng túng xua tay như thể bị khen ngợi quá đà: ``E-Em nghĩ văn phòng cần chút gì đó nhẹ nhàng, thư giãn hơn\dots\ nên em mới chọn cái này\dots''

Kuon gật gù đồng tình với lời giải thích ngắn gọn đó, giọng đầy tán thành: ``Ừ, hợp lý lắm. Mấy cái này để trên kệ hay bàn tiếp khách đều rất đẹp.''

Trong khi mọi người đang bận rộn khoe những món đồ vừa mua, Ryuuhou vẫn đứng lặng ở một góc phòng, tay ôm chặt chiếc túi của mình với vẻ mặt đầy bí ẩn. Cô bé liếc nhìn mọi người một lượt rồi mỉm cười, như đang giữ một bí mật thú vị chỉ riêng mình biết.

Kuon nhìn em gái với ánh mắt tò mò: ``Đến lượt em rồi đấy, Ru-chan. Em mua cái gì thế?''

Ryuuhou cười tinh nghịch, giấu chiếc túi ra sau lưng: ``Để sau hẵng hay, Onii-san. Đợi mọi người xong xuôi hết em sẽ khoe\~{}''

Polka cười gùn gò, khoanh hai tay trước ngực với vẻ đắc ý của kẻ sắp gây ra một vụ sốc bất ngờ: ``Đúng đúng, đến lượt em rồi đây nhỉ!''

Cô thình lình lôi một cái túi vải cồng kềnh từ sau lưng tung ra, rồi úp ngược nó xuống sàn nhà. Một chùm thú nhồi bông hình hề với nụ cười méo xệch, đôi mắt lồi hẳn ra lập tức lăn lộn khắp nền, trông chẳng khác nào đám sinh vật vừa trốn thoát khỏi cơn ác mộng của ai đó.

Không khí trong phòng lập tức thấm đẫm sự rờn rợn, chẳng ai cất lên được tiếng nào. Mọi người im lặng nghệt mặt, không biết phản ứng thế nào trước cô nàng Cáo xiếc vẫn đang hí hửng như thể vừa cướp được kho báu của vua nào đem về.

Chẳng mấy chốc, Miko, Mio và Korone cùng lúc kêu lên với tiếng hét lẫn lộn kinh ngạc và sợ hãi:

Nàng Vu nữ lắp bắp, giọng run rẩy như vừa thấy ma: ``\textbf{N-N-N-Ne!? Cái gì mà lắm thế này ne!?}''

Nàng Sói đen cũng không kém phần hoảng hốt: ``\textbf{Cái quái gì thế này!?}''

Nàng Khuyển lùi vội một bước, tiếng kêu lạc đi vì sợ: ``\textbf{Uwwooh!? Trông chúng nó như sắp nhảy bổ vào mặt người ta vậy!}''

Kuon vẫn đứng yên khoanh tay nhìn đống thú bông, rồi từ từ hỏi với giọng bối rối như đang cố giải một bài toán vô lý: ``Ờm\dots\ Sao em lại mua mấy thứ này về hả?''

Polka cười lớn, đặt tay lên hông đầy kiêu hãnh, giọng pha chút triết lý nghệ thuật: ``Thì nhóm idol cũng phải có phong cách riêng chứ! Đây chính là đặc trưng của Polka!''

Kuon chỉ biết cười trừ bất lực, giọng giễu cợt: ``Anh thấy em đem chúng đi trang trí nhà ma còn hợp hơn đấy\dots\ Ai bước vào thấy cái này là khóc thét!''

Mio ôm đầu rên rỉ, giọng than thở của kẻ đã hết hy vọng: ``Không chịu nổi nữa rồi\dots\ Bây giờ trong phòng vừa có cục đá kỳ lạ, quả lựu đạn phát sáng, lại thêm đám hề quái dị này\dots\ may mà còn có mấy lọ hoa của Lamy để cân bằng lại chút!''

``Sao em nói vậy, cái đèn của Botan cũng ổn mà,'' Kuon lên tiếng, cố tìm một điểm sáng trong mớ hỗn loạn.

``\textbf{Không hề ổn chút nào!!}'' nàng Sói đen phản đối lập tức, giọng đầy quyết liệt như đang chống lại một điều phi lý.

Cậu Tổng thở dài, liếc qua Miko và Polka rồi ra lệnh với giọng muốn lập lại trật tự: ``Được rồi hai đứa, mau cất mấy thứ kinh dị của mình vào góc đi, đừng để ai vào văn phòng tưởng đây là nơi phong ấn yêu quái.''

``Ehhh, vậy là em không được đặt cục đá này ở giữa phòng hả?'' nàng Vu nữ rên rỉ tiếc nuối, giọng như đứa trẻ bị tịch thu đồ chơi yêu thích.

``\textbf{Dĩ nhiên là không được!!}'' cả Mio và Kuon cùng hét lên, giọng đồng thanh như một bản hợp xướng của sự kiên quyết không thể thay đổi.

Polka thở dài đành gom đám hề nhồi bông của mình lại, xếp gọn vào một góc phòng với vẻ mặt như bị ép phải chia tay những người bạn thân. Còn Miko thì khác, cô nàng vẫn nhất quyết giữ nguyên cục đá Moai ngay giữa phòng, đôi mắt sáng rực đầy quyết tâm.

``Miko-san vẫn chưa chịu dời cục đá đi ạ?'' Ryuuhou lên tiếng, nhìn chủ nhân của nó đang tự sướng chụp ảnh với cục đá như một nghệ sĩ đang tôn vinh tác phẩm đỉnh cao của mình.

``Ne ne, nó là điểm nhấn độc đáo mà! Biết đâu ai đó lại thích nó chứ,'' cô đáp, đôi mắt ánh lên sự quả quyết bất chấp mọi lời chê, như người đang bảo vệ một ý tưởng nghệ thuật đầy tranh cãi.

``Chị có biết ai lại thích mấy cục đá mặt người như này đâu!'' Mio đáp lại giọng dằn mặt, như một lời cảnh báo cuối cùng.

Vừa khi cuộc tranh cãi vẫn đang đi đến đâu không biết, một giọng nói nhẹ nhàng nhưng đầy vững tin vang lên từ phía cuối phòng. Ryuuhou tiến lại với nụ cười rạng rỡ, trên tay cầm một chiếc hộp nhỏ nhắn đã được quấn gói cẩn thận.

``Onii-san, đến lượt em khoe đồ rồi nè!'' cô bé thốt lên, giọng tràn đầy nôn nao như một đứa trẻ đang nóng lòng trổ tài thành tích học tập của mình trước cả nhà.

Cả nhóm lập tức xoay người nhìn về phía cô bé mắt hai màu với ánh mắt tò mò dồn dập, trong khi cậu Tổng nhõng một bên mày lên: ``Ồ? Cái gì mà bí mật quá vậy, em mua gì thế?''

Ryuuhou từ từ mở nắp hộp ra. Bên trong là một bộ khung ảnh gỗ được chạm khắc tỉ mỉ, cộng thêm một chiếc kệ treo tường nhỏ gọn và một chùm đèn dây xinh xắn. Những hoa văn trên gỗ tuy đơn giản nhưng rất thanh thoát, mang lại cảm giác ấm áp và có chút nghệ thuật tự nhiên.

``Em nghĩ văn phòng của chúng ta cần một chút không gian sáng tạo, và em học được điều này từ những lần tự dựng sân khấu cho ban nhạc underground em tham gia,'' cô giải thích với giọng của một người đã có kinh nghiệm thực tế. ``Mấy cái khung ảnh này có thể để ảnh của các chị, kệ treo thì đặt mấy chậu cây nhỏ hay sách, còn đèn dây thì treo trên tường để không gian thêm ấm cúng.''

Mio nhìn những món đồ Ryuuhou vừa mang ra mà đôi mắt sáng rực, giọng cô như vừa gặp được vị cứu tinh: ``Cuối cùng cũng có thêm một người suy nghĩ bình thường! Em quá giỏi rồi, Ryuuhou-chan!''

Kuon cũng nhìn em gái với nét mặt tự hào, giọng đầy hài lòng: ``Phải công nhận món này hợp lý thật. Em học được mấy cái này từ đâu vậy?''

Ryuuhou cười khúc khích, giọng như đang chia sẻ những mẹo nghề nhỏ: ``Hồi em làm trưởng nhóm, bọn em hay phải tự tay trang trí toàn bộ không gian biểu diễn. Chỉ cần vài chi tiết nhỏ thôi là có thể thay đổi cả không khí của cả căn phòng đấy, Onii-san!''

Lamy cũng gật đầu lia lịa, giọng đầy đồng tình: ``Ý tưởng của Ryuuhou-chan hay thật đấy. Đơn giản mà lại tạo cảm giác dễ chịu lắm.''

Polka nhìn những món đồ của cô bé, rồi liếc qua đám thú nhồi bông hề của mình, cô bé lẩm bẩm nhỏ với vẻ mặt hơi có chút ghen tỵ: ``Cũng đẹp thật\dots\ nhưng phong cách của Polka vẫn là nhất mà!''

Cùng lúc đó, Korone - vốn dĩ từ nãy giờ chỉ đứng từ xa quan sát - bỗng giơ lên một bình xịt màu vàng với vẻ mặt vô cùng tự hào, giọng như sắp giới thiệu một phát minh thay đổi thế giới: ``Em cũng mua được một món rất hữu dụng này nè!''

Cậu Tổng và nàng Sói đen cùng dồn ánh nhìn vào cái bình trên tay nàng Khuyển. Đó là một bình xịt diệt côn trùng, trên nhãn có vẽ hình một con gián bị gạch chéo, sản phẩm của hãng của nàng Vịt.

``\dots Nhưng đây có phải là đồ trang trí đâu em?'' cậu nghiêng đầu, vẻ mặt khó hiểu đến mức không dám tin vào mắt mình.

``Mà\dots\ chẳng lẽ \emph{mấy con đó} thật sự có bò quanh trong văn phòng này sao!?'' nàng Sói đen rùng mình, nổi cả da gà, tay cô ôm chặt lấy cánh tay mình như đang tự trấn an trước một mối đe dọa vô hình.

Korone gật đầu một cách vô cùng nghiêm chỉnh, giọng như một chiến binh đang chuẩn bị ra trận: ``Nếu không có thì tốt quá, nhưng mà nếu có thì tôi đã sẵn sàng xử lý chúng rồi!''

Nàng Yêu tinh tuyết mặt tái nhợt, lùi ra xa một bước với vẻ mặt hoảng sợ: ``Ơ\dots\ nhưng mà\dots\ sao chị lại mua cái đó về vậy?''

Nàng Cáo Sa mạc vội bịt chặt hai tai, lắc đầu nguầy nguậy không muốn nghe thêm một chữ nào: ``Đừng nói nữa! Đừng nhắc đến cái tên đó! Chỉ cần nghĩ tới thôi là em đã thấy lạnh cả sống lưng rồi!''

Nàng Sư tử trắng khoanh tay trước ngực, cười nhẹ với giọng bình thản như không có gì to tát: ``Tốt là có chuẩn bị, nhưng nếu thật có thì chắc chúng chạy mất từ lâu rồi, trước khi em kịp xịt đâu.''

``Đúng vậy. Thêm nữa, anh nghĩ bọn chúng cũng chả đủ can đảm để bén mảng tới đây đâu, kể từ ngày chúng ta dọn về tới giờ.''

``Anh chỉ nghĩ thôi hả!? Sao không chắc chắn 100\% đi!?'' Polka hoảng loạn, hai tay vẫn giữ chặt tai, giọng đầy bất an.

Lamy tuyết ôm lấy đầu, giọng run rẩy như vừa nghe thấy tin dữ: ``Hôm nay em về sớm một chút vậy\dots''

Korone chớp mắt nhìn cả nhóm phản ứng quá đà, rồi lẩm bẩm một cách ngây thơ: ``Ơ, vậy ra em là người bình tĩnh nhất ở đây à?''

Cậu Tổng xoa trán, quyết định không tranh cãi thêm về vấn đề đó nữa. Nhưng ngay lúc đó, nàng Vu nữ bỗng gọi cả nhóm lại với giọng hào hứng, như một đứa trẻ vừa khám phá ra trò chơi mới: ``Ne, nhìn này, nhìn này mọi người ơi!''

Nàng Sói đen nhăn mặt quay sang, giọng mệt mỏi như đã chịu đựng đủ mọi chuyện lạ lùng trong ngày: ``Lại có chuyện gì nữa đây?''

Miko hớn hở tạo dáng nâng niu cục đá Moai nhỏ trên tay, cười phá lên đến mức nước mắt chảy dài, giọng vui mừng như vừa thắng được giải thưởng lớn: ``Moai của em cầm trên tay này!!''

\img{7-13l}

Nàng Sói đen và cậu Tổng nhìn cảnh đó mà chẳng thể nào cảm thấy có gì vui vẻ. Cả hai cùng đồng thanh la lên với giọng bất lực tột độ.

Mio ra lệnh với giọng đã hết kiên nhẫn: ``\textbf{Đem cất nó đi ngay!}''

Kuon cũng không kém phần gay gắt, giọng đanh thép như lời phán quyết cuối cùng: ``\textbf{Trông nó chẳng hợp với căn phòng chút nào!}''

Ryuuhou đứng bên cạnh nhìn mọi người cãi nhau, cô bé vừa thấy buồn cười vừa thích thú, rồi thì thầm nhỏ vào tai anh trai: ``Onii-san, em nghĩ Miko-san sẽ không bao giờ chịu dời cục đá đó đâu. Nhưng mà ít nhất văn phòng cũng có thêm vài món trang trí đẹp từ mọi người rồi.''

Cậu Tổng thở dài nhìn em gái với vẻ mặt bất lực: ``Chỉ là mấy món đồ đẹp ấy lại rơi vào tay những người không bình thường thôi.''

``Nhưng mà như thế mới vui chứ anh!''

Trong lúc cả nhóm đang còn đau đầu loay hoay với mớ đồ độc nhất vô nhị kia, một giọng nói êm dịu xen lẫn chút tò mò vang lên từ phía sau: ``Ủa? Mấy đứa đang bận làm gì thế?''

Miko giật mình quay người liền, đôi mắt lập tức sáng rực như vừa gặp được vị cứu tinh: ``Sora-chan!?''

Kuon cũng ngước mắt nhìn về phía cửa, gật đầu với vẻ điềm nhiên như thường lệ: ``Em đến sớm thế ha.''

Nhưng trước khi kịp có ai phản ứng thêm, Mio đột nhiên mặt tái nhợt, nuốt nước bọt ừng ực đầy lo lắng, tựa như một con thú nhỏ vừa chợt nhận ra kẻ săn mồi đang ở ngay trước mặt. Cậu Tổng cũng hơi lùi chân lại một chút, thì thầm trong hơi thở: ``C-Chết thật rồi\dots''

Ryuuhou đứng nép bên cạnh anh, cũng cảm nhận được bầu không khí vừa chuyển biến lạ thường. Cô bé thì thầm hỏi anh trai: ``Onii-san, sao Sora-san đến mà mọi người có vẻ sợ hãi thế nhỉ?''

Cậu Tổng cũng nhỏ giọng đáp lại: ``Không phải sợ chị ấy đâu, mà là sợ cái ánh mắt đó thôi\dots'' Cậu vừa kịp nói vậy thì nàng Thần tượng bước vào, nghiêng đầu với nụ cười dịu dàng như thường\dots\ thế nhưng ánh mắt cô lại lạnh lẽo sắc bén đến đáng sợ, tựa một lưỡi dao bén ngót bị che giấu sau lớp vải nhung mềm mại: ``Cái gì mà chết rồi? Nói to ra cho chị nghe nào.''

\img{7-13m}

Nàng Sói đen rùng mình, ôm lấy đầu kêu nhỏ với giọng run rẩy như một đứa trẻ sắp bị mắng: ``Híc\~{}!''

Nàng Vu nữ cũng đỏ bừng cả mặt, sợ hãi đến mức nước mắt cứ sắp trào ra, giọng cô lạc đi như đang cầu cứu: ``K-Không có gì đâu\dots!!''

Riêng Korone thì chỉ đứng ngẩn ngơ, hoàn toàn không hiểu nổi sao bầu không khí lại thay đổi đột ngột đến thế. Cô nghiêng đầu sang một bên, gương mặt trông y hệt một chú chó đang cố gắng hiểu một mệnh lệnh quá phức tạp.

Botan lùi chân lại một bước, đưa tay che miệng rồi thì thầm vào tai cô bạn Yêu tinh tuyết, giọng như đang bình luận về một tình huống đầy rẫy nguy hiểm: ``\hl{Mình cứ cảm thấy vừa chứng kiến một cảnh tượng gì đó rất đáng ngại đúng không?}''

Lamy cũng e dè đáp lại, giọng đầy bỡ ngỡ như vừa khám phá ra một điều gì đó hoàn toàn mới lạ: ``\hl{Đây có phải là\dots{} cái uy quyền của một người tiền bối không nhỉ\dots?}''

Polka thì run lẩy bẩy, tay túm chặt lấy cánh tay cô bạn Sư tử trắng, giọng thì thầm lo lắng: ``\hl{Hay là mình nên cười lớn một cái cho không khí bớt căng thẳng nhỉ\dots?}''

Cậu Tổng thì bình tĩnh hơn một chút, chỉ lẩm bẩm với giọng của một người đã quá quen với những tình huống dở khóc dở cười này: ``Mới đầu xuân thôi mà sao cảnh tượng gì như sắp vào hè thế\dots''

Ryuuhou vẫn đứng nép sau lưng anh trai, vừa tò mò vừa e ngại nhìn nàng Thần tượng, rồi cũng nhỏ giọng hỏi: ``Onii-san, em thấy Sora-san trông cũng dễ thương mà, sao mọi người lại sợ chị ấy đến thế?''

Cậu Tổng cúi đầu xuống nhỏ tiếng đáp: ``Em chưa từng thấy em ấy nổi giận đâu. Tin anh đi, em chắc chắn sẽ không bao giờ muốn chứng kiến cảnh đó.''

Bỗng dưng, gương mặt của Sora giãn ra hoàn toàn, cô bật cười lớn như một đứa trẻ vừa thành công trong một trò đùa tinh nghịch, chắp tay xin lỗi: ``Chị đùa thôi mà! Mấy đứa đang định sửa sang, trang trí lại cái văn phòng này đúng không?''

Cả bảy người vẫn chưa kịp hoàn hồn sau cú sốc vừa rồi thì nàng Thần tượng đã liền tay trao cho mỗi người một vật nhỏ xinh, khiến ai cũng phải nhìn chằm chằm vào bàn tay mình với vẻ bối rối, tất cả ánh mắt đều đổ dồn vào những món đồ lấp lánh vừa có được.

\img{7-13n}

``Cái này là gì vậy ạ?'' nàng Khuyển ngây thơ hỏi, đôi mắt mở to tròn đầy tò mò.

Nàng Thần tượng mỉm cười tự hào, đặt tay lên hông với giọng của một người đang trao phần thưởng xứng đáng: ``Đó là tem `Cảm ơn vì đã giúp đỡ!' đấy!''

Cậu Tổng chớp mắt vài lần để xử lý tình hình, rồi lặng lẽ lên tiếng với giọng đang cố gắng hết sức để giữ bình tĩnh: ``\dots Nhưng mà anh có cần cái này đâu chứ! Bọn anh có phải là mấy đứa trẻ mẫu giáo đâu mà cũng phải nhận tem khen thưởng!''

Cô vẫn tươi cười rạng rỡ, tạo dáng đáng yêu như một idol chuyên nghiệp trên sân khấu: ``Chỉ có những bé ngoan mới được nhận nó thôi mà, đúng không nào?''

Ryuuhou nhìn cái tem trên lòng bàn tay mình, mỉm cười vui vẻ: ``Em thấy cái tem này xinh lắm! Cảm ơn Sora-san nhiều ạ!''

Kuon vốn định phản bác thêm nhưng rồi cũng đành thôi, chỉ có thể thở dài bất lực chấp nhận số phận, giọng cậu như một tiếng thở dài của kẻ vừa đầu hàng. Cậu quay sang nhìn nàng Vu nữ, người đang cầm chặt cái tem trên tay mà người cứ run lên cầm cập, dường như vui sướng đến mức sắp bật khóc.

``Miko-chan, sao mà em phải rung lắc thế? Lần đầu tiên trong đời được nhận tem thưởng à?''

Mio cũng nhìn cái tem trong tay mình, rồi thở dài não nề với giọng của một người đã chính thức từ bỏ mọi hy vọng: ``Nói thật thì, mấy cái tem này\dots''

Miko liền tiếp lời, giọng hào hứng như đang khen ngợi một món ăn cực kỳ ngon miệng: ``\dots nhìn thật sự rất là đáng yêu, thấy muốn sưu tầm lắm ne.''

Cậu Tổng buông ra một tiếng thở dài cuối cùng, hoàn toàn không biết phải phản ứng thế nào nữa, tay cứ liên tục xoa lên trán với vẻ bất lực tột độ.

Polka thì lật qua lật lại cái tem trong tay, tặc lưỡi với giọng đánh giá như đang xem xét một món đồ có giá trị: ``Ít nhất cũng có chút quà an ủi cho bao công sức chúng ta bỏ ra hôm nay.''

Lamy ngập ngừng nhìn đàn chị, rồi lại ngó xuống cái tem trong tay, hỏi một cách băn khoăn: ``Nhưng mà\dots\ chị còn có tem vàng nào không ạ?''

Sora nghiêng đầu tò mò: ``Tem vàng là cái gì thế?''

Botan nhếch mép cười nhẹ, lên tiếng giải thích cho trò đùa của bạn mình: ``Ý Lamy là hỏi là có món quà gì to tát, hoành tráng hơn không đấy.''

Nàng Thần tượng liền giả vờ bĩu môi ra vẻ bất bình: ``Ôi trời đất ơi, giờ đến cả cái tem `Cảm ơn bạn đã giúp!' cũng bị chê bai nữa hả?''

Nàng Cáo Sa mạc vội vàng xua tay cuống quýt, giọng hoảng hốt như vừa làm sai chuyện gì lớn: ``Không không không, em nhận, em nhận mà! Cảm ơn Sora-senpai rất nhiều!''

Ryuuhou đứng bên cạnh cũng cười khúc khích trông rất thích thú, lên tiếng hỏi: ``Sora-san này, chị còn có thêm tem nào cho em không ạ? Em thấy cái tem này đẹp lắm, muốn sưu tầm thêm!''

Nàng Thần tượng mỉm cười thân thiện, liền đưa thêm một cái tem cho Ryuuhou: ``Của em đây nhé!''

\il

Trong khi cả nhóm vẫn còn bàn tán về đống tem ``Cảm ơn bạn đã giúp!'' của Sora, ở phía xa, một ánh mắt lấp lánh nhưng đầy phức tạp đang dõi theo họ.

Ryuuhou với đôi mắt tinh tường, bỗng nhiên nhíu mày nhìn về phía cánh cửa, cô bé thì thầm với anh trai: ``Onii-san, hình như lại có ai đó đang nhìn tụi mình kìa.''

Cậu Tổng khẽ liếc theo hướng em gái chỉ, nhưng không để lộ sự chú ý: ``Anh biết rồi. Cứ để yên.''

Người đó đứng nấp sau cánh cửa, má phồng lên một cách đáng yêu, đôi tay siết nhẹ gấu áo. Không phải vì cô ghen tị với mấy cái tem trẻ con kia; cô chả cần mấy thứ đó! Nhưng mà\dots

``Mọi người vui vẻ quá ha\dots''

Ánh mắt cô dừng lại không phải trên nàng Sói đen, nàng Vu nữ hay nàng Khuyển, mà là trên ba người đồng đội của mình - nàng Cáo Sa mạc, nàng Yêu tinh tuyết và nàng Sư tử trắng - đang đứng ngay bên cạnh cậu Tổng, thoải mái cười nói với cậu.

Cô nàng bĩu môi, bất mãn, giọng cô như một tiếng thì thầm đầy ghen tị: ``Mấy người đó là đội của mình mà, sao lại ở bên cạnh cậu ta thế chứ? Mình cũng muốn một phần trong đó!''

Cô em gái của cậu Tổng như thể có thể cảm nhận được sự hiện diện của người lạ, cô bé lén lút nhìn về phía cánh cửa và thì thầm với anh trai: ``Onii-san, em thấy người đó cứ đứng đó mãi. Chị ấy có vẻ muốn tham gia nhưng không dám.''

Kuon khẽ gật đầu: ``Để em ấy tự quyết định.''

Cô nhìn chằm chằm vào cảnh tượng trước mặt, rồi đột nhiên siết chặt tay thành nắm đấm, ánh mắt lóe lên như vừa quyết định điều gì đó.

``Được rồi! Mình cũng phải thể hiện mới được!''

Và thế là, nàng idol tóc vàng lao ra khỏi chỗ nấp với một quyết tâm mãnh liệt, sẵn sàng tham gia vào cuộc trò chuyện\dots\ dù cô chưa nghĩ ra là mình sẽ nói gì. Nhưng trái tim cô thì đang đập mạnh với những dự cảm về một cuộc gặp gỡ sắp diễn ra.

%==========================================TRUYỆN PHỤ=========================================
\omk

Nhận ra cái tem ``Cảm ơn vì đã giúp!'' mà Sora - hay đúng hơn là công ty - trao cho những ai hoàn thành tốt nhiệm vụ, Korone và Miko đã quyết định thực hiện mọi việc để có được những chiếc tem ấy. Dù về mặt lý thuyết chúng chẳng có giá trị gì, y hệt những tấm tem dán thư thông thường, nhưng với hai cô gái, đó là cả một niềm tự hào lớn lao.

Thậm chí, hai người còn dấn thân vào những hành động kỳ quặc nhất. Một trong số đó chính là trở thành những chiến binh chuyên diệt gián.

``Mio-senpai, chị vừa nói là trong văn phòng có nhiều con đó đúng không?'' nàng Khuyển tò mò hỏi, cái đuôi của nàng vẫy nhẹ nhàng, khuôn mặt hệt một thợ săn đang săn tìm con mồi.

Ryuuhou đứng gần đó chớp mắt với vẻ khó hiểu: ``Ở văn phòng này có nhiều con như vậy đó sao?''

Korone quay sang em gái, giọng đầy bí ẩn: ``Ưm! Mấy con có râu dài, chạy lướt như bay, và\dots'' nàng chưa kịp nói hết thì Mio đã lắc đầu ngắt lời.

``Không đâu, chị chỉ mới nói là `chẳng lẽ văn phòng này bị \emph{bọn chúng} xâm nhập rồi sao' thôi,'' cô đáp, tay vẫn đang bê một chậu cây nhỏ, vẻ mặt rõ ràng là không muốn nhắc đến loài vật ấy.

Thấy vậy, Kuon thở dài với giọng của một người đã quá quen với những nỗi sợ vô lý: ``\emph{Gián} thì cứ nói thẳng ra thôi. Sao phải sợ ghê vậy??''

Vừa nghe cái tên ấy, cả ba người lập tức run lên cầm cập. Korone và Miko vừa chạy khắp phòng vừa la hét như những đứa trẻ bị dọa ma, còn Mio thì nhảy thẳng lên trên bàn với khuôn mặt đầy kinh hoàng. Cả ba cùng đồng thanh hét lên: ``\textbf{Đừng có nhắc đến tên nó!}''

Ryuuhou thấy các chị phản ứng quá khích thì đưa tay che miệng cười khúc khích: ``Trông các chị buồn cười quá! Có gì đâu mà phải sợ thế chứ? Chỉ là mấy con hay chạy lông nhông mỗi khi hè về thôi mà. Nhà em toàn gặp mấy con ấy suốt!''

Kuon nhìn em gái với vẻ bất lực, rồi quay sang ba cô gái đang hoảng loạn. Cậu chẹp miệng với giọng như đang chứng minh một điều hiển nhiên: ``Xời, mấy con đó thì có gì mà phải sợ. Đây này, có một con đấy.''

Nói là làm, cậu chỉ tay về phía một con gián đang bò thong thả trên sàn nhà, chẳng hề có chút cảnh giác gì. Ba cô gái kia lại tiếp tục la hét, mặt tái xanh trước sinh vật gớm ghiếc, như thể vừa nhìn thấy một con quái vật từ hành tinh khác.

``\textbf{Giết nó đi! Giết mau đi!}'' nàng Sói đen hét lớn, chỉ vào con gián như thể đó là một mối đe dọa lớn nhất, giọng đầy tuyệt vọng.

``\textbf{Rõ!}''

Ngay lập tức, nàng Khuyển và nàng Vu nữ lao vào: người cầm bình xịt côn trùng vừa mua xịt thẳng vào con vật, kèm theo một luồng khói trắng; người thì vung cây gậy trừ tà đuổi theo, động tác hệt một hiệp sĩ đang chiến đấu với rồng.

Chẳng mấy chốc, dưới tác dụng của thuốc xịt và mấy cái quật mạnh của cây gậy, con gián đã nằm im lìm, chịu một cái kết chẳng mấy vẻ vang.

Ryuuhou nhìn cảnh đó vừa buồn cười vừa thán phục: ``Các chị làm nó chết thảm quá! Bọn em cùng lắm chỉ cần dùng một cái dép đập là được mà!''

Ba cô gái vừa rồi hoảng loạn là thế, ba thành viên còn lại cũng không khá hơn mấy - trừ Botan.

Lamy từ nãy giờ đứng sát sau hai người bạn, đột nhiên bấu chặt lấy tay đồng đội, mắt nhắm nghiền lại, giọng run rẩy như một đứa trẻ vừa xem xong phim kinh dị: ``C-Chuyện\dots\ chuyện vừa rồi\dots\ là thật có đúng không\dots?''

Botan nhìn cái xác bất động trên sàn, thở hắt ra một tiếng với giọng của một người đang bình luận về một trận chiến: ``Ít nhất thì nó cũng chết rồi. Nhưng cảnh này trông cứ như một trận chiến khốc liệt thật sự\dots''

Polka, thay vì sợ hãi, lại bỗng dưng tỏ ra hứng thú. Cô nghiêng đầu nhìn cái bình xịt trong tay Korone, lẩm bẩm như đang tính toán một kế hoạch lớn: ``Hay đấy\dots\ Vậy có nghĩa là cứ tiêu diệt được mấy con này là mình sẽ có tem `Cảm ơn vì đã giúp' hả? Hừm\dots''

Ý tưởng ấy làm nàng Yêu tinh tuyết sợ xanh mặt, giật mình quay sang cô bạn với giọng đầy lo lắng: ``B-Bà định làm cái gì thế!?''

Cô nàng chủ xiếc cười xòa, vẫy tay nhẹ nhàng như thể chỉ đang đùa: ``Không không, đùa thôi! Tôi đâu có ham mấy cái tem đó đến thế. Tôi cũng sợ gián chết khiếp mà.''

Ryuuhou nhìn Polka với vẻ tò mò: ``Polka-san nói thế em cũng thấy sợ theo luôn đấy! Nhưng mà diệt gián thì có được tem thật không ạ?''

Kuon nhìn em gái với vẻ bất lực, lắc đầu: ``Em đừng có học theo mấy chị ấy nhé.''

Mãi khi thấy con gián không còn cử động gì nữa, Mio mới hết run, nhưng giọng vẫn còn hơi lạc đi, như một người vừa thoát khỏi cơn ác mộng: ``Nó\dots\ chết thật chưa?''

Miko và Korone đồng thanh đáp, giọng đầy tự hào như vừa hoàn thành một nhiệm vụ lớn: ``Chết rồi chị ạ!''

Cậu Tổng chỉ đứng đó cười ha hả, như đang thưởng thức một vở hài kịch hay nhất đời, trong khi Lamy vẫn còn ôm chặt cánh tay cô bạn Sư tử, như thể chỉ cần buông ra là sẽ có một con gián khác liền xuất hiện.

``Anh còn đứng đó cười cái gì nữa!? Mau vứt nó đi đi!'' Mio la lên từ trên bàn, nước mắt đã bắt đầu rơi, giọng cầu cứu đầy tuyệt vọng.

``Rồi rồi, dọn ngay đây.'' Cậu lấy một cái hót rác và cây chổi, quét xác con gián lên rồi vứt ra ngoài cửa sổ với động tác dứt khoát, như thể đã làm việc này cả đời.

Sau đó, nàng SÓi đen mới chịu bước xuống sàn, còn Lamy thì vẫn đứng sát bên cô bạn lực lưỡng như một con mèo vừa bị dọa sợ đang nấp sau người lớn, đôi mắt còn long lanh vì sợ hãi.

Trong khi cả nhóm vừa thở phào nhẹ nhõm sau khi đuổi đi kẻ xâm nhập, bộ ba thành viên của NePoLaBo chỉ đứng nhìn mà chẳng dám tham gia vào trận chiến đầy kinh hoàng kia. Có lẽ chỉ trừ Botan, người luôn có thể cười đùa với mọi chuyện.

Lamy nắm chặt hai bàn tay trước ngực, mặt trắng bệch như vừa chứng kiến một cơn ác mộng, giọng thì thầm đầy sợ hãi: ``K-Khủng khiếp quá\dots\ may mà nó không bò về phía tôi\dots.''

Botan khoanh tay, lắc đầu ngao ngán với giọng của một người đã chứng kiến quá nhiều sự ồn ào: ``Vậy mà mới nãy còn nói là không có gì phải sợ\dots''

Polka rụt cổ lại, rón rén lại gần Mio, mắt nhìn quanh phòng đầy nghi ngờ như một thám tử đang tìm manh mối: ``Ê, chắc chỉ có một con thôi đúng không? Hay nó còn có em, có cháu gì ở đây nữa!?''

Nghe vậy, cô bạn Yêu tinh tuyết hoảng loạn bấu chặt lấy cánh tay Botan, giọng đầy sợ hãi: ``Đ-Đừng có nói thế mà! Tôi không muốn biết đâu!''

Cả nhóm tưởng đã có thể quên đi nỗi kinh hoàng, thì đột nhiên Kuon lên tiếng với nụ cười tinh nghịch, như kẻ chuẩn bị kể chuyện ma: ``Mấy đứa có muốn anh kể cho nghe một câu chuyện về gián không?''

``\textbf{Im đi!!}''

Mio lập tức bịt miệng cậu, trong khi Miko và Korone vừa hoảng loạn vừa tức tối vung tay đập nhẹ vào lưng Kuon, như một đàn chim đang mổ vào kẻ gây rối.

``Người ta vừa hết sợ mà chị lại định kể chuyện kinh dị hả!? Không có nhân tính đâu ne!!'' nàng Vu nữ vung cây gậy trừ tà của mình vào người Kuon, giọng đầy phẫn nộ.

``Không phải người bình thường đâu, chắc chắn luôn!'' nàng Sói đen cũng hoảng loạn, vừa khóc thét vừa đấm đá túi bụi vào người cậu Tổng, giọng như một đứa trẻ đang khóc nức nở.

Còn nàng Khuyển thì đơn giản là hùa theo vì thấy vui, đôi mắt sáng rực như một đứa trẻ đang chơi trò đuổi bắt.

Ryuuhou đứng bên cạnh, nhìn anh trai bị mọi người tấn công vừa buồn cười vừa thương hại: ``Onii-san, em nghĩ lần này anh đã chọc các chị quá trớn rồi đấy.''

Bị cả ba cô gái cùng lúc tấn công, Kuon chỉ có thể cười chịu trận, giơ hai tay lên đầu hàng: ``Rồi rồi, anh im miệng. Không trêu nữa.''

Ryuuhou nhìn cảnh tượng với vẻ mặt đầy thích thú, rồi thì thầm với nàng Sư tử trắng đang đứng cạnh: ``Botan-san, em thấy mọi người lúc nào cũng vui vẻ thế này à?''

Nàng Sư tử trắng mỉm cười, đáp lại: ``Ừ, em sẽ quen thôi. Đây là \hll\ mà.''

Thế nhưng, dù nỗi sợ đã tạm lắng xuống, một vấn đề khác lại dấy lên. Nàng Cáo Sa mạc liếc nhìn nàng Sói đen đầy nghi ngờ, giọng cô như một thám tử đang thẩm vấn: ``Này này, nhưng chị vẫn chưa trả lời câu hỏi của em đấy! Chắc chắn là chỉ có một con thôi đúng không?''

Không một ai dám chắc chắn về điều đó. Và thế là, cả đám lại bắt đầu đảo mắt khắp phòng, đề phòng một cuộc tấn công bất ngờ từ ``bọn đồng loại'' của kẻ vừa bị xử lý và đuổi đi. Ánh mắt của mọi người quét sạch mọi ngóc ngách, đầy cảnh giác như những lính canh đang chờ đợi kẻ thù xông tới.

Ryuuhou đứng ở một góc, vừa nhìn quanh vừa tỏ ra tò mò thích thú: ``Mấy chị sợ gián đến thế à? Em thấy nó cũng chỉ là một loại côn trùng nhỏ thôi mà.''

Lamy quay sang cô bé với giọng đầy bất lực: ``Em chưa từng thấy chúng bay đâu! Đợi lúc chúng vỗ cánh bay nhảy, em sẽ hiểu cảm giác đó thôi!''

Từ đó, cả nhóm bắt đầu họp chiến lược, những cái đầu tụm lại với nhau như một ban tham mưu đang lên kế hoạch tác chiến.

Lamy vẫn còn ôm chặt lấy Botan, mắt vẫn còn đảo quanh như thể con gián vừa bị vứt đi chỉ là kẻ đi đầu cho một đội quân đang ẩn nấp: ``Chúng\dots\ chúng còn có thể chui ra từ đâu nữa không ạ?''

Cô bạn Sư tử trắng khoanh tay, ngắm nghía xung quanh với vẻ mặt của một vị tướng đang thăm dò chiến trường: ``Nếu nói về lý thuyết, bọn chúng thích những nơi tối và ẩm ướt. Có thể là đằng sau tủ, dưới gầm bàn, hay\dots\ là cả trên trần nhà nữa.''

``Trần\dots\ trần nhà á!?''

Lamy lập tức run cầm cập, ngước cổ lên nhìn trần nhà với vẻ hoảng sợ, lo rằng bất cứ lúc nào cũng có một thứ gì đó rơi xuống đầu mình, như thể cái trần nhà bỗng dưng trở thành mối đe dọa khủng khiếp nhất.

Polka đứng bên cạnh cũng nhìn theo hướng cô đang nhìn, nhưng thay vì sợ hãi, cô bỗng nảy ra một ý tưởng với ánh mắt rạng rỡ của một nhà phát minh vừa tìm ra giải pháp: ``À này, hay là chúng ta sắp xếp lại phòng một chút để dụ bọn chúng ra rồi dẹp sạch luôn?''

Ryuuhou nghe vậy liền vỗ tay thích thú: ``Ý hay đấy! Chúng ta có thể dùng đống đồ trang trí để bịt hết những chỗ tối tăm, thế là gián chẳng còn nơi nào để trốn nữa!''

Kuon nhìn nàng Cáo xiếc một lúc, rồi gật đầu đồng tình với giọng của người đã quá quen với những kế hoạch quái gở của mọi người: ``Nghe cũng được. Tận dụng đợt sửa sang này để làm chúng mất chỗ ẩn náu luôn.''

Mio nghe vậy liền giật mình, giọng đầy lo lắng: ``Khoan đã! Đừng bảo là mọi người định thực sự làm bẫy dụ chúng ra thật đấy!? Chị không muốn chứng kiến cảnh bọn chúng ồ ạt kéo đến như hồi chị dọn nhà cho Kuon-senpai hồi Tết năm ngoái đâu!''

Lamy vội vàng gật đầu lia lịa, giọng như đang cầu xin: ``Đúng đấy! Mio-senpai, em chịu không nổi cảnh đó đâu!''

Botan vỗ nhẹ vai cô bạn để trấn an, giọng đầy chắc chắn: ``Đừng lo. Nếu chúng có xuất hiện, cứ giao cho đội đặc nhiệm `Xịt là chết' xử lý hết.'' Cái tên đội đặc nhiệm mà Botan vừa nghĩ ra lập tức làm Miko và Korone phấn khích. Cả hai cùng giơ cao bình xịt lên với vẻ mặt hào hứng của những chiến binh sẵn sàng ra trận.

``Chúng ta sẽ dẹp sạch hết lũ xấu xa đó! Này Korone, em có nghĩ bọn chúng còn có `trùm cuối' không?'' Miko nói, vừa nói vừa lôi thêm một bình xịt nữa từ túi đồ của Korone ra, chuẩn bị cho một trận đánh boss lớn.

Korone cũng nghiêm túc gật đầu, tỏ ra hứng thú với giọng của một game thủ đang chờ đợi thử thách cuối cùng: ``Chắc chắn rồi! Biết đâu lại có một con `đại boss' to hơn đang ẩn nấp ở đâu đó!''

Ryuuhou cười khúc khích, cũng cầm lấy một bình xịt nhỏ: ``Em cũng tham gia cùng mọi người nhé! Dù em không sợ gián, nhưng thấy mọi người nhiệt tình quá nên em cũng muốn góp vui!''

Kuon vừa ngắm nhìn bộ ba đã sẵn sàng lao vào nhiệm vụ sinh tử kia, vừa thở dài bất lực, khoanh tay lại với nụ cười không thể nào tin được: ``Trời đất, mấy đứa làm nghiêm túc quá, cứ như đây là game chiến dịch thật sự vậy\dots''

Liền sau đó, không khí phòng lập tức sôi động khi ba thành viên của đội đặc nhiệm ``Xịt là chết'' bắt đầu dồn dập chuẩn bị vũ khí. Miko đầu tiên lôi ra từ túi đồ hai bình xịt côn trùng còn nguyên hộp, phấn khích xé tung bao bì đến nỗi miếng băng keo bắn ra xa, rồi trao một bình cho Korone: ``Của em đây! Cái này công suất xịt siêu mạnh, bảo đảm gián nào cũng phải nằm im!'' 

Korone ôm chầm lấy bình xịt như thể đang cầm một khẩu súng trường, rồi lại lôi thêm từ đâu ra cây gậy trừ tà quen thuộc, gõ nhẹ xuống sàn nhà với tiếng ``thình thịch'' đầy uy lực. 

Ryuuhou cũng không chịu thua kém, cô bé chạy nhanh nhặt thêm một cái bình xịt nhỏ còn sót lại, rồi nghiêm chỉnh xếp tất cả vũ khí vào một cái túi vải, vừa làm vừa lẩm bẩm như một quân nhân đang kiểm đếm đạn dược: ``Còn có thêm hai cái dép nhựa nữa, để nếu có con nào chạy nhanh thì em có thể đập kịp!'' Cả ba xung nhau kiểm tra vũ khí đến mức cười nói rộn cả phòng, trông hệt như một đội chiến binh sẵn sàng ra trận.

Polka hào hứng vỗ tay liên tục, giọng tràn đầy phấn khích như một MC đang dẫn chương trình thực tế: ``Ôi hay quá! Thế là hôm nay chúng ta vừa hoàn thành nhiệm vụ trang trí, vừa có thêm giải trí với trận đại chiến diệt gián sao!? Nghe quá thú vị rồi!''

Lamy thì mặt mày tái xanh đi sau khi nghe vậy, giọng run rẩy như sắp khóc: ``Không, đâu có thú vị gì đâu, tôi thấy chỉ có đáng sợ thôi\dots''

Mio cũng lắc đầu ngao ngán, cố gắng hít một hơi sâu để lấy lại bình tĩnh, rồi quay sang cậu Tổng với giọng vẫn còn đầy lo lắng: ``Nhưng mà\dots\ em cũng sợ lắm, sợ là lại có con khác chui ra từ đâu đó, nên là\dots''

Cô không thể nói hết câu vì sợ đến mức lạnh cả gốc gáy, tựa một con thỏ nhỏ đang đối mặt với nguy hiểm rình rập. Kuon liền tiếp lời thay cho cô, giọng đầy khích lệ như một vị chỉ huy động viên binh lính: ``Ý của Mio-chan là muốn mấy đứa giúp ngăn chặn bọn chúng đấy. Cứ làm tốt đi, biết đâu Sora-chan lại thưởng thêm tem cho cả đội.''

Miko và Korone lập tức hiểu ý cậu Tổng, hào hứng cầm lấy bình xịt mà phun khắp mọi ngóc ngách trong phòng để tiễn đưa hết lũ côn trùng đáng ghét, tựa như những người lính đang quét sạch hoàn toàn kẻ thù ra khỏi căn cứ.

\img{7-13o}

Ryuuhou cũng nhiệt tình tham gia, cô bé phun thuốc vào những góc khuất khó thấy rồi kêu lên: ``Miko-san, Korone-san, em phụ trách xịt bên này nhé!''

``Anh nghĩ mấy đứa xịt hơi bị nhiều quá rồi đấy\dots'' Kuon lẩm bẩm, khi cả căn phòng đang phủ đầy một màn sương trắng bay lượn, tạo nên cảnh tượng hư ảo chẳng khác nào những bộ phim khoa học viễn tưởng trên ti vi.

Ngược lại, Mio - người cũng đang cầm một bình xịt khác - lại thúc giục mọi người với giọng vẫn còn hoảng loạn: ``Cứ xịt nhiều vào! Chỉ có vậy bọn chúng mới không còn đường nào mà trốn thoát được!''

Miko, Korone và Ryuuhou cùng đồng thanh đáp lại như những người lính nhận lệnh, giọng nói đầy nghiêm chỉnh: ``\textbf{Rõ!}''

Cậu Tổng chỉ biết chẹp miệng mặc kệ bọn họ phun thuốc khắp văn phòng, trong khi cậu bắt đầu dọn dẹp và bày biện những đồ trang trí. Từ những chú thú nhồi bông đến khối đá lớn mà Miko đã mang về - tự lúc nào nó cũng đã được thu nhỏ thành phiên bản mini, xinh xắn. Ngay cả bốn hàng cây cảnh cậu mua cũng được mọi người tận dụng để phục vụ cho công việc chung.

Polka giúp sắp xếp lại những món đồ sao cho hài hòa hơn, còn Botan thì phụ Kuon di chuyển những vật nặng. Lamy, dù trong lòng vẫn còn e ngại, cũng cố gắng góp sức bằng cách đặt những chậu cây nhỏ lên bàn, với đôi bàn tay vẫn còn run rẩy vì sợ.

Sau khi xong xuôi việc phun thuốc, Ryuuhou chạy lại giúp anh trai treo những khung ảnh và dây đèn nhấp nháy mà cô bé đã mua sắm: ``Onii-san, em treo mấy cái đèn này lên tường nhé? Em nghĩ treo lên trông sẽ ấm cúng lắm!''

Kuon nhìn em gái với nụ cười tự hào, gật đầu đồng ý: ``Ừ, em cứ làm đi.''

\ct

Chỉ mười lăm phút sau, mọi công việc đều hoàn thành. Bảy người cùng đứng nhìn thành quả của mình, ai nấy đều tràn đầy sự hài lòng, những gương mặt rạng rỡ chiếu sáng dưới ánh đèn mới lắp.

``Anh nghĩ thế này là ổn thỏa rồi nhỉ?'' Kuon đứng chống hông, không giấu được vẻ tự hào về những gì cả nhóm đã làm được.

``Chúng ta đã dùng hết tất cả đồ đạc rồi, thế này là quá tuyệt luôn!'' Mio hào hứng thốt lên, giọng cô tràn đầy sự nhẹ nhõm sau bao nhiêu căng thẳng.

``Nhiệm vụ xịt phòng cũng đã hoàn thành 100\%!'' Korone nói sau khi phun thêm một chút vào những khe hở cuối cùng, rồi giơ tay chào kiểu nhà binh với vẻ mặt đầy kiêu hãnh.

``Cả căn phòng bây giờ trông tươi sáng hơn hẳn rồi ne!'' Miko đặt cây chổi lông (cũng chính là cây gậy trừ tà của mình) xuống, vừa lau mồ hôi trán vừa trông mệt mỏi nhưng vô cùng hài lòng.

Ryuuhou nhìn quanh căn phòng với đôi mắt long lanh, giọng cô ngập tràn niềm vui: ``Phòng đẹp quá đi! Mấy cái dây đèn và khung ảnh này trông ấm cúng muốn chết!''

Polka cũng ngắm nghía khắp phòng một lượt rồi gật gù, giọng cô như một nhà phê bình đang đánh giá một tác phẩm nghệ thuật: ``Không tệ chút nào! Chúng ta vừa đánh bại được lũ gián, lại vừa hoàn thành xong việc trang trí văn phòng. Một công đôi việc luôn!''

Botan khoanh tay cười, giọng cô đầy vẻ thỏa mãn: ``Chà, đúng là một ngày vất vả thật sự.''

Lamy thì khẽ thở phào nhẹ nhõm, cuối cùng cũng cảm thấy yên tâm như vừa thoát ra khỏi một cơn ác mộng: ``Cuối cùng cũng xong rồi\dots\ Chắc không còn con nào chui ra nữa đâu nhỉ?''

Ngay sau đó, nàng Vu nữ và nàng Chó kéo nàng Sói đen sang phòng quản lý kế bên để đòi Sora đưa những chiếc tem ``Cảm ơn vì đã giúp!''. Kuon cũng mỉm cười đi theo sau bọn họ. Ba cô gái của NePoLaBo cũng hí hửng đuổi kịp, tiếng cười nói vang vọng khắp cả hành lang văn phòng.

Ryuuhou cũng vội vàng chạy theo, cô bé giơ tay gọi lớn: ``Chờ em với! Em cũng muốn xin thêm tem nữa!''

Như vậy, một mùa đông lạnh giá đã trôi qua, để lại văn phòng với một diện mạo hoàn toàn mới: rộng rãi hơn, tươi tắn hơn và thoáng đãng hơn bao giờ hết. Những món đồ trang trí được sắp xếp hài hòa, những chậu cây xanh mướt tỏa ra sức sống mãnh liệt, và những dãy đèn dây ấm áp treo trên tường, toát lên sự gần gũi.

Một mùa xuân mới đang đến gần, và những bước chuẩn bị cuối cùng để đón những ngày ấm áp cũng đã hoàn thiện. Căn phòng đã sẵn sàng chào đón một mùa mới tràn đầy hy vọng.

\il

\pov{Hikawa Kuon}

Sau cả ngày dài vừa ``chiến đấu'' với lũ gián vừa dọn dẹp văn phòng, ai trong chúng tôi cũng có chút mệt nhoài, nhưng tâm trạng thì vô cùng thoải mái, như thể vừa mới trút bỏ được một gánh nặng bấy lâu đè trên vai. Tất cả mọi người đều đã về từ lâu, còn con bé Ryuuhou thì lon ton chạy về nhà trước để chuẩn bị cơm nước.

Tôi rải bước dọc con đường vắng lặng cùng ba cô bạn đồng hành: Botan, Lamy và Polka. Mặt trời chiều đang dần lặn sau dãy nhà cao tầng, nhuộm cả bầu trời một gam màu cam nhạt ấm áp như vắt sữa. 

``Thế là xong xuôi hết rồi nhỉ\dots'' Cáo Sa mạc vừa nói vừa duỗi người thỏa mái, giọng cô như một tiếng thở dài mãn nguyện, ``Tôi cứ tưởng vụ này sẽ kéo dài cả ngày cơ chứ.''

Yêu tinh tuyết gật đầu lia lịa, tay vẫn siết chặt chiếc khăn quàng hình gấu mình, khuôn mặt còn nguyên vẻ hoảng hốt vừa trải qua một cuộc phiêu lưu đầy rẫy nguy hiểm: ``Ừm\dots\ may mắn là mọi chuyện cũng êm xuôi rồi.''

Nàng Sư tử trắng khoanh tay trước ngực, nở một nụ cười hài lòng, giọng cô như một lời đồng tình khẳng định thành quả chung của cả nhóm: ``Đúng thật. Văn phòng bây giờ trông khác xa ngày hôm qua luôn.''

Tôi chỉ mỉm cười khẽ, im lặng bước đi cùng các cô ấy, để những dòng suy nghĩ của mình trôi dạt theo làn gió chiều se lạnh. Một mùa xuân nữa sắp gõ cửa, chẳng biết những điều gì đang chờ đợi chúng tôi ở phía trước? Những kế hoạch mới đang ấp ủ, những dự định chưa kịp vạch ra, và có lẽ cả những rắc rối khó lường mà chúng tôi chẳng thể nào ngờ tới.

Mọi thứ đáng lẽ sẽ kết thúc một cách êm đẹp như vậy, nếu như không có\dots

``Ể? Sao mọi người về sớm thế?''

\dots người ấy xuất hiện.

Chúng tôi đồng loạt dừng chân khi nghe giọng nói ấy vang lên. Phía trước, một cô gái với mái tóc vàng óng ánh, phía dưới có điểm thêm vài lọn tóc màu hồng, đang đứng đó. Cô ấy trông có vẻ ngạc nhiên khi tình cờ gặp chúng tôi, như thể vừa bắt gặp một điều thú vị ngoài dự đoán.

Lamy là người đầu tiên lên tiếng, giọng cô đầy bất ngờ: ``Ơ? Bà\dots?''

Botan - người mà tôi thường đùa gọi là ``người chồng'' của cô bạn - cũng hơi nhướng mày, giọng cô lộ rõ sự quan tâm xen lẫn ngạc nhiên: ``Không phải bà vẫn còn đang ốm ở nhà sao?''

Cô gái ấy - thành viên cuối cùng của nhóm NePoLaBo - nở một nụ cười thân thiện, giọng nói tuy còn hơi yếu ớt sau khi khỏi bệnh nhưng vẫn tràn đầy sức sống, như một ngọn đèn vừa được thắp sáng trở lại sau khi tắt lụn: ``Mình khỏi bệnh rồi! Cuối cùng cũng thoát khỏi cái cơn cảm cúm khốn khổ ấy!

Polka chớp mắt vài lần trước khi bật cười vui mừng cho cô bạn, giọng cô như một tiếng reo hò: ``Tuyệt quá! Cả tháng nay không thấy bóng dáng bà đâu, bọn tôi còn tưởng bà bị biến khỏi mặt đất luôn rồi cơ.''

``Đúng vậy thật,'' Lamy cũng gật gù đồng tình, giọng cô tràn đầy sự nhẹ nhõm, ``Mấy hôm trước tôi còn định nhắn tin hỏi thăm sức khỏe nữa, mà sợ bà đang cần thời gian nghỉ ngơi nên thôi\dots''

Tất cả chúng tôi đều vui mừng khi thấy cô ấy đã khỏe mạnh trở lại, dù giọng nói vẫn còn hơi khàn đi. Thế nhưng cô ấy lại bỗng dưng phồng má lên một cách đáng yêu, khoanh tay trước ngực với giọng đầy hờn dỗi của một đứa trẻ bị bỏ rơi: ``Hứ! Mấy người toàn đi gặp Kuon-senpai mà chẳng bao giờ rủ tôi lấy một lần!

Nàng Sư tử trắng cười khúc khích, liếc mắt sang tôi với vẻ thích thú: ``Thì tại bà đang ốm phải nằm dưỡng sức mà làm sao rủ được?''

``Cho nên tôi mới ghen tị lắm!'' cô ấy giậm chân nhẹ xuống vỉa hè, khuôn mặt lộ rõ vẻ ấm ức y hệt một đứa trẻ vừa bị tịch thu món đồ chơi yêu thích, ``Cứ như thể tôi là người ngoài cuộc ấy! Mọi người được gặp Kuon-senpai suốt ngày, còn tôi thì bị gạt hẳn ra ngoài\dots''

Tôi nhìn cô ấy trong giây lát, trong lòng dấy lên những cảm xúc hỗn loạn mà chẳng thể nào gọi tên được. Trong ánh mắt của cô ấy, ngoài sự hờn dỗi ra còn có chút gì đó khó tả\dots\ Tôi chẳng dám chắc đó là ánh mắt háo hức mong chờ được gặp tôi, nhưng cũng không phải là không có chút gì đó như thế.

Hay có lẽ là bởi cô ấy coi tôi là một thứ gì đó đặc biệt, thật sự xứng đáng? Cũng không hẳn. Chắc là tôi lại đang suy diễn lung tung thôi, vẫn như mọi lần.

Cô ấy bỗng quay sang phía tôi, chống nạnh với khuôn mặt vẫn còn phồng má kia, giọng cô như một lời thách thức đầy quyết tâm: ``Kuon-senpai! Anh đừng nghĩ là anh có thể trốn tránh em mãi nhé!''

Tôi nhún vai, đáp lại bằng giọng trêu chọc quen thuộc của một người đã quá quen với những lời đe dọa vô hại này: ``Anh có trốn đâu đâu. Tại em không lên văn phòng gặp anh đấy chứ.''

Cô ấy bĩu môi, đôi mắt bỗng sáng bừng lên với một tia tinh nghịch: ``Hừm, dù sao đi nữa\dots\ Em cũng định lên văn phòng sớm thôi! Lúc đó đừng có mà cố tình lảng tránh em đấy!''

Nói xong, cô ấy nở một nụ cười rạng rỡ - vẫn là nụ cười tràn đầy năng lượng mà tôi vẫn nhớ, một nụ cười như thể vừa tìm được mục tiêu mới để phấn đấu trong cuộc đời.

Tôi chỉ cười nhẹ, cảm thấy có chút gì đó ấm áp len lỏi trong lòng: ``Được rồi. Anh sẽ đợi em.''

Cô ấy mỉm cười lần nữa, rồi vẫy tay chào tạm biệt tất cả mọi người trước khi bước đi. Bóng dáng nhỏ nhắn của cô dần dần khuất sau con phố đông đúc, hòa vào dòng người đang hối hả với những công việc cuối ngày.

Polka nghiêng đầu nghi ngờ, khoanh tay trước ngực như một thám tử đang phân tích một nghi phạm, giọng cô thì thầm: ``Này, có ai thấy bà ấy hơi\dots\ khác lạ không?''

Botan  nhếch mép suy nghĩ, giọng cô như một nhà phân tích đang đưa ra giả thuyết hợp lý nhất: ``Hừm, có lẽ là do ốm đau lâu ngày nên tâm trạng thay đổi thôi chứ sao?''

Lamy đặt một tay lên ngực, thì thầm với giọng khó hiểu của một người vừa cảm nhận được một điều gì đó mơ hồ: ``Nhưng tôi thấy\dots\ cậu ấy có vẻ vui hơn bình thường đấy.'' 

Cả ba cô gái đồng loạt quay sang phía tôi, như thể đang chờ đợi một câu trả lời từ tôi, những ánh mắt họ đầy tò mò.

Tôi không nói gì, chỉ lặng lẽ nhìn theo hướng cô ấy vừa rời đi.

Một ngày không xa, chúng tôi sẽ gặp lại nhau. Và có lẽ, đó sẽ lại là một trò mèo gì đó làm cho tôi phát điên đây\dots

\dots M███████ N███ ạ.

\end{document}